% !TeX program = pdflatex
%
% Paper A of River's odd-zeta programme.
% Working title: "Lucas congruences for the second solutions of Apery-like recurrences".
%
% EVERY statement in this file is traceable to one of the source memos listed in
% Section "Provenance of the statements" (sec-data.tex).  Citations are restricted to
% references that carry an arXiv identifier in those memos; the remaining references
% were verified separately against their sources (see work/BIBLIO_VERIFY.md).
%
\documentclass[11pt,reqno]{amsart}

\usepackage[T1]{fontenc}
\usepackage[utf8]{inputenc}
\usepackage{amsmath,amssymb,amsthm}
\usepackage{mathtools}
\usepackage{booktabs}
\usepackage{array}
\usepackage{longtable}
\usepackage{enumitem}
\usepackage[margin=1.15in]{geometry}
\usepackage{microtype}
\usepackage[hidelinks]{hyperref}

% ---------------------------------------------------------------- theorem styles
\theoremstyle{plain}
\newtheorem{theorem}{Theorem}[section]
\newtheorem{proposition}[theorem]{Proposition}
\newtheorem{lemma}[theorem]{Lemma}
\newtheorem{corollary}[theorem]{Corollary}
\newtheorem{conjecture}[theorem]{Conjecture}

\theoremstyle{definition}
\newtheorem{definition}[theorem]{Definition}
\newtheorem{example}[theorem]{Example}
\newtheorem{problem}[theorem]{Problem}

\theoremstyle{remark}
\newtheorem{remark}[theorem]{Remark}
\newtheorem{observation}[theorem]{Observation}

% ---------------------------------------------------------------- notation
\newcommand{\Z}{\mathbb{Z}}
\newcommand{\Q}{\mathbb{Q}}
\newcommand{\N}{\mathbb{N}}
\newcommand{\Zp}{\Z_{(p)}}
\newcommand{\vp}{v_p}
\newcommand{\bin}[2]{\binom{#1}{#2}}
\newcommand{\HH}[2]{H^{(#1)}_{#2}}
\newcommand{\Kchi}[2]{K^{(#1)}_{\chi}(#2)}
\newcommand{\lcm}{\operatorname{lcm}}
\DeclareMathOperator{\ct}{ct}

% evidence-class markers.  These are load-bearing: they are never upgraded.
\newcommand{\Proved}{\textup{\textsf{[PROVED]}}}
\newcommand{\Verified}{\textup{\textsf{[VERIFIED]}}}
\newcommand{\VerifiedR}[1]{\textup{\textsf{[VERIFIED: #1]}}}
\newcommand{\Open}{\textup{\textsf{[OPEN]}}}

\numberwithin{equation}{section}

% ---------------------------------------------------------------- front matter
\begin{document}

\title[Lucas congruences for second solutions]
      {$\chi$-twisted Lucas congruences for the second solutions\\ of Ap\'ery-like recurrences}

%% FIXED [MINOR]: author placeholder was being typeset. Left empty, as in the companion
%% FIXED drafts, to be filled in before circulation.
\author{}

\date{\today}

\subjclass[2020]{11A07, 11B65, 11S80, 33F10}
\keywords{Ap\'ery numbers, Lucas congruences, Dwork congruences, second solution,
sporadic sequences, $p$-adic Ap\'ery limits, creative telescoping}

\begin{abstract}
Let $(A(n))_{n\ge 0}$ be one of the fifteen sporadic Ap\'ery-like sequences, and let
$(B(n))_{n\ge 0}$ be the second solution of its three-term recurrence, normalised by
$B(0)=0$, $B(1)=1$. The first row satisfies a Lucas congruence,
$A(ap+r)\equiv A(a)A(r)\pmod p$; for the second row no congruence of Lucas type appears
in the literature. We prove that the second row descends along base-$p$ digits as well,
subject to two corrections: a normalising factor $p^{w}$, where $w$ is the total degree of
the sequence's harmonic weight, and a quadratic Dirichlet character $\chi$. Concretely, for the
classical Ap\'ery pair for $\zeta(3)$ we prove
\[
  p^{3}\,b_{ap+r}\;\equiv\;b_{a}\,a_{r}\pmod p
  \qquad(p\ge 5,\;1\le a<p,\;0\le r<p),
\]
by an elementary self-contained argument, and we then prove a general theorem: if the
second solution of a sporadic pair admits a harmonic-monomial weight of total degree $w$
%% FIXED [MINOR x2]: (i) Theorem \ref{thm:LB} gives chi(p)^e, not chi(p); (ii) w and chi are
%% FIXED defined here through the Apery limit, which three of the fifteen families do not have.
satisfying five explicit hypotheses, then $p^{w}B(ap+r)\equiv\chi(p)^{e}B(a)A(r)\pmod p$, with
$e\in\{0,1\}$ read off the weight. The
character is not decorative. The untwisted law fails completely for seven of the fifteen
families, at exactly the primes where $\chi(p)=-1$, and the single factor $\chi(p)$ --
which we trace to one line of Frobenius descent for a $\chi$-twisted harmonic letter --
repairs it. For the twelve families that have an archimedean Ap\'ery limit, $w$ and $\chi$
are its weight and character; for the other three both are read off the harmonic weight
directly, and the descent holds regardless. Exact computation
over all primes $5\le p\le 103$ supports a master form,
$\vp\bigl(p^{w}B_nA_q-\chi(p)B_qA_n\bigr)\ge w$ with $q=\lfloor n/p\rfloor$, with floor
exactly $w$ in every cell tested; we state it as a conjecture and never as a theorem.
Three of the fifteen sequences have no archimedean Ap\'ery limit at all, and satisfy the
$p$-adic descent regardless. Along the way we prove a new closed form for the Ap\'ery
numerators,
$b_n=\sum_k\binom nk^2\binom{n+k}k^2\bigl(2\HH{3}{n}-\HH{3}{k}\bigr)$,
which we have not located in the literature and which reduces the weight-three congruence
to three lines of the general theorem. The general theorem, both first-row Lucas
congruences and two instances of the second-row congruence are formalised in Lean~4,
sorry-free. In the language of Beukers--Vlasenko's Dwork crystals our congruence concerns
the non-unit diagonal entry of the Frobenius matrix, complementary to their
Conjecture~7.5 about the unit root; the character is invisible in unit-root formulations
and has to be put in by hand.
\end{abstract}

\maketitle

{\small\tableofcontents}

\section{Introduction}
\label{sec:intro}

\subsection{The first row}
\label{ssec:firstrow}

Let
\begin{equation}\label{eq:apery-a}
  a_n \;=\; \sum_{k=0}^{n}\bin nk^2\bin{n+k}k^2
  \;=\;1,\;5,\;73,\;1445,\;33001,\;819005,\;21460825,\;\dots
\end{equation}
be Ap\'ery's numbers for $\zeta(3)$ (OEIS \texttt{A005259}). They satisfy a congruence of
Lucas type: for every prime $p$, every $a\ge 0$ and every $r$ with $0\le r<p$,
\begin{equation}\label{eq:arow}
  a_{ap+r}\;\equiv\;a_a\,a_r \pmod p ,
\end{equation}
and consequently $a_n\equiv\prod_i a_{n_i}\pmod p$ for the base-$p$ expansion
$n=\sum_i n_ip^i$. This is classical \cite{Gessel1982}, and it is by now known
to be a structural, not accidental, feature of the whole \emph{sporadic} class:
Malik and Straub \cite[Thm.~3.1]{MalikStraub} prove \eqref{eq:arow} for every one of the
fifteen sporadic Ap\'ery-like sequences in Zagier's classification, by McIntosh's
constant-term method, and Straub \cite{Straub2301} sharpens it to a congruence modulo
$p^2$ for all fifteen. Beyond Lucas congruences there is a supercongruence:
$a_{p^rm}\equiv a_{p^{r-1}m}\pmod{p^{3r}}$ for $p\ge 5$
\cite[Thm.~3.2]{Straub1401}, a statement whose history runs through Beukers' conjecture
and Coster's proof, as recorded in the survey of Osburn and Sahu \cite{OsburnSahu}. On the
crystalline side, Beukers and Vlasenko's Dwork-crystal papers
\cite{BVone,BVtwo,BVthree} and Mellit--Vlasenko \cite{MellitVlasenko} put the same
first-row congruences into a Frobenius framework, where they appear as statements about
the \emph{unit root} $F(t)/F(t^\sigma)$ of the holomorphic period.

In short: the arithmetic of the integral solution of an Ap\'ery-like recurrence is
well-trodden ground.

\subsection{The second row, and the gap}
\label{ssec:gap}

Ap\'ery's proof of the irrationality of $\zeta(3)$ uses a \emph{second} solution of the
same recurrence,
\begin{equation}\label{eq:apery-b}
  b_n \;=\; \sum_{k=0}^n \bin nk^2\bin{n+k}k^2
     \Bigl(\HH{3}{n} \;+\; \sum_{m=1}^{k}
       \frac{(-1)^{m-1}}{2m^3\bin nm\bin{n+m}m}\Bigr),
  \qquad \HH{r}{x}=\sum_{m=1}^{x}\frac1{m^r},
\end{equation}
with $b_n/a_n\to\zeta(3)$; the first values are
$0,\,6,\,\tfrac{351}{4},\,\tfrac{62531}{36},\,\tfrac{11424695}{288},\dots$
(OEIS \texttt{A059415} for the numerators). The pair $(a_n,b_n)$ is not integral: $b_n$
has denominators, bounded by $d_n^3b_n\in\Z$ with $d_n=\lcm(1,\dots,n)$.

For this second row, nothing of the shape \eqref{eq:arow} appears in the literature. We
surveyed the relevant technology in detail, working only from fetched sources, and the
picture is uniform:

\begin{itemize}[leftmargin=1.4em]
\item \emph{Lucas congruences for constant-term sequences} (Malik--Straub
  \cite{MalikStraub}, Henningsen--Straub \cite{HenningsenStraub}, Mellit--Vlasenko
  \cite{MellitVlasenko}, Straub \cite{Straub1401,Straub2301}) treat the integral solution
  only.
\item \emph{Dwork crystals} \cite{BVone,BVtwo,BVthree} construct the logarithmic second
  solution $F(t)\log t+G(t)$, whose coefficients are exactly our $B(n)$; they explain the
  $p^w$ normalisation, since the Cartier matrix on the rank-two quotient is
  $\Lambda=Y\Lambda_0(Y^\sigma)^{-1}$ with
  $\Lambda_0=\begin{psmallmatrix}\alpha_0&\alpha_1\\0&p\alpha_0\end{psmallmatrix}$, so
  that the second-solution Frobenius eigenvalue carries exactly one extra factor of $p$
  per Hodge step \cite[Prop.~7.7]{BVthree}; and they prove $p$-integrality of the mirror
  map $q(t)=t\exp(G/F)$ \cite[Cor.~7.11]{BVthree}. What they do \emph{not} prove is any
  congruence modulo $p$ or $p^s$ for the second solution: the corresponding deepening is
  their Conjecture~7.5, explicitly open, and stated for the \emph{unit root} at that.
\item \emph{Mirror-map integrality} (Delaygue--Rivoal--Roques \cite{DRR}, Delaygue
  \cite{Delaygue}, \cite{MirrorNote}) and \emph{strong Frobenius structures}
  (Vargas-Montoya \cite{VMcanon,VMmum}) control the denominators of the second-solution
  ratio, and in \cite{VMmum} prove Lucas congruences --- but explicitly only for the
  \emph{holomorphic} solution.
\item \emph{$p$-adic zeta values} sit in the Frobenius matrix entries of a MUM operator
  (Beukers--Vlasenko \cite{BVfrob}), which is the conceptual reason a weight-$w$ second
  solution should descend with a factor $p^w$ --- but that is a statement about a $p$-adic
  \emph{limit}, proved for simplicial and hyperoctahedral Calabi--Yau families, and the
  classical Ap\'ery operator is not among them.
\item The one genuine precedent is Gessel's companion congruence, in the form given by
  Beukers \cite[Thm.~1.4]{Beukers2211}: with
  $A'_k=\sum_{m=0}^k\bin km^2\bin{m+k}m^2\bigl(\tfrac1{m+1}+\dots+\tfrac1k\bigr)$ one has
  $A_{kp+\ell}\equiv A_\ell A_k+p\,A'_\ell\,k\,A_k \pmod{p^2}$. This is a
  \emph{weight-one} harmonic companion, and the depth stops at $p^2$. (We have confirmed
  the modulus $p^2$ and the definition of $A'$ against the source; the placement of the
  indices in the $p$-term should be checked against the published version before it is
  relied upon.)
\end{itemize}

So the second row is a genuine gap: weight one is known modulo $p^2$; weight three, the
weight of $\zeta(3)$ and of Ap\'ery's own companion, is not known at all.

\subsection{Results}
\label{ssec:results}

\subsubsection*{(A) The weight-three theorem}
Our first result is elementary and self-contained.

\begin{theorem}[weight-three descent; Theorem~\ref{thm:brow} below]
\label{thm:intro-brow}
Let $p\ge 5$ be prime, let $1\le a<p$ and $0\le r<p$, and put $n=ap+r$. Then $b_a$ and
$p^3b_n$ lie in $\Zp$ and
\[
   p^3\,b_{ap+r}\;\equiv\;b_a\,a_r \pmod p .
\]
\end{theorem}

The proof occupies Section~\ref{sec:weight3}; it needs Lucas' theorem, Kummer's carry
criterion, and nothing else. The three steps that a referee found compressed in the first
write-up --- the borrow count in $\bin nm$, the product-completion of the summation
region, and the interchange defining the tail sums --- are here isolated as
Lemma~\ref{lem:borrow}, Lemma~\ref{lem:fibre} and Lemma~\ref{lem:interchange}
respectively, and proved in full.

\subsubsection*{(B) The master depth law}
Numerically the congruence of Theorem~\ref{thm:intro-brow} is the mod-$p$ shadow of a
sharper and more symmetric statement which we can verify but not prove. Writing
$q=\lfloor n/p\rfloor$, the quantity $\vp(p^3b_na_q-b_qa_n)$ is $\ge 3$, with minimum
exactly $3$, for every prime $5\le p\le 31$ and every $n\le 320$ --- including the
multi-digit range and including $p=5$. Section~\ref{sec:depth} states this as
Conjecture~\ref{conj:master3} and records the exact sweep parameters, the fine structure
of the sporadic boosts above the floor, and two false statements that the data kills. The
descent depth equals the motivic weight of the Ap\'ery limit; this is what makes the
programme's name for it, \emph{the master depth law}, appropriate.

\subsubsection*{(C) The general $\chi$-twisted theorem}
The classification of Ap\'ery-like recurrences produces fifteen sporadic pairs
$(A(n),B(n))$. Running the same experiment on all fifteen produces the paper's main
discovery: the naive law is \emph{false} for seven of them, and the failure is governed by
a quadratic Dirichlet character.

\begin{conjecture}[master form; Conjecture~\ref{conj:master} below]
\label{conj:intro-master}
For each sporadic pair, with $w$ the weight of its Ap\'ery limit, $\chi$ the quadratic
character attached to it, $p\ge 5$ and $q=\lfloor n/p\rfloor$,
\[
  \vp\bigl(p^{w}B_nA_q-\chi(p)B_qA_n\bigr)\;\ge\;w .
\]
\end{conjecture}

\noindent
\VerifiedR{zero failures} for all fifteen pairs, all primes $5\le p\le 103$, all $n\le400$,
and $n\le 1200$ for $p=5,7,11$ (five base-$p$ digits), in exact rational arithmetic; the
floor is exactly $w$ in every single (sequence, prime) cell.

The mod-$p$, single-digit form of Conjecture~\ref{conj:intro-master} we can prove, in a
generality that covers four of the fifteen families outright:

\begin{theorem}[Theorem~\ref{thm:LB} below]
\label{thm:intro-LB}
Let $p\ge5$ and let $B(n)=\sum_kS(n,k)w(n,k)$, where $S$ is a product of binomials in
linear forms and $w$ is a $\Q$-combination of weight-$w$ monomials in harmonic letters
$\Kchi{r}{x}=\sum_{m\le x}\chi(m)m^{-r}$. Under the five hypotheses
\textup{(H1)--(H5)} of Section~\ref{sec:general},
\[
  p^{w}\,B(ap+r)\;\equiv\;\chi(p)^{e}\,B(a)\,A(r)\pmod p
  \qquad(1\le a<p,\;0\le r<p),
\]
where $e\in\{0,1\}$ is the common number of $\chi$-letters per monomial.
\end{theorem}

The mechanism is a two-line lemma (Lemma~\ref{lem:K}): for a completely multiplicative
$\chi$,
\[
   p^r\Kchi ry \;=\; \chi(p)\,\Kchi r{\lfloor y/p\rfloor}\;+\;p^r\cdot(\text{$p$-integral}),
\]
because the $p$-divisible indices $m=jp$ contribute $\chi(jp)(jp)^{-r}
=\chi(p)p^{-r}\chi(j)j^{-r}$. \emph{The factor $\chi(p)$ in the congruence is exactly the
$\chi(p)$ coming out of the $p$-divisible layer of the harmonic weight.} Correspondingly,
$\chi$ is the character of the sequence's Ap\'ery limit: the limit is a rational multiple
of $\zeta(w)$ when $\chi$ is trivial, and of $L(\chi_D,w)$ when $\chi=\chi_D$. The
correlation is $12$ out of $12$ on the sequences whose limit exists.

\subsubsection*{(D) A minimal closed form}
The general theorem needs a harmonic-monomial weight, and finding one for Ap\'ery's $b_n$
produced a formula we have not been able to locate anywhere:
\begin{equation}\label{eq:minimal-intro}
   b_n\;=\;\sum_{k=0}^n\bin nk^2\bin{n+k}k^2\bigl(2\HH3n-\HH3k\bigr).
\end{equation}
Section~\ref{sec:minimal} proves \eqref{eq:minimal-intro} by six explicit
Zeilberger/Gosper certificates, and proves two further closed forms of the same kind, for
Franel's numbers and for $\sum_k\bin nk^4$. Substituted into Theorem~\ref{thm:intro-LB},
\eqref{eq:minimal-intro} reduces Theorem~\ref{thm:intro-brow} to three lines. It also
removes the hypothesis $p\ge5$ for that instance, since the coefficients $2,-1$ are
integers.

\subsubsection*{(E) Sharpness, positioning, formalisation}
Section~\ref{sec:sharpness} documents the failure of the master law at $p=2,3$, and
locates the obstruction where the proof says it should be: in the denominators of the
weight's coefficients. Section~\ref{sec:discussion} positions
Conjecture~\ref{conj:intro-master} against Beukers--Vlasenko's Conjecture~7.5: the two are
statements about the two diagonal entries of one $2\times2$ Frobenius matrix, neither
implies the other, and the character is invisible in the unit-root formulation.
Section~\ref{sec:lean} reports what is machine-verified: Theorem~\ref{thm:intro-LB},
Lemma~\ref{lem:K}, the two first-row Lucas congruences, and two instances of the
second-row congruence are formalised in Lean~4, sorry-free, on the three standard Mathlib
axioms.

\subsection{Evidence classes}
\label{ssec:evidence}

This paper reports a mixture of proofs and of exact finite computations, and it labels
them. We use exactly three markers, and we never upgrade one to another:

\begin{description}[leftmargin=2.2em,style=nextline]
\item[\Proved] a complete proof is written out here.
\item[\Verified] an exact finite computation, always with its range stated. Exact rational
  or modular arithmetic, no floating point. This is evidence and never proof, however
  large the range.
\item[\Open] neither.
\end{description}

Statements carrying only \Verified\ are called \emph{conjectures}, and are stated in
\texttt{conjecture} environments, with the sweep parameters attached. Where a computation
was run twice, by independent implementations, we say so.

\subsection{Notation}
\label{ssec:notation}

Throughout, $p$ is a prime, $\vp$ is the $p$-adic valuation on $\Q$ normalised by
$\vp(p)=1$, and $\Zp=\{x\in\Q:\vp(x)\ge0\}$ is the local ring; a congruence
$x\equiv y\pmod{p^s}$ between rationals means $\vp(x-y)\ge s$. We write $d_n=\lcm(1,\dots,n)$,
$\HH rx=\sum_{m=1}^x m^{-r}$ (so $\HH r0=0$), and for a completely multiplicative
arithmetic function $\chi$,
\[
   \Kchi ry\;=\;\sum_{m=1}^{y}\frac{\chi(m)}{m^{r}} ,
\]
which for $\chi\equiv1$ is $\HH ry$. For a sporadic pair we write $A(n)$ for the integral
solution and $B(n)$ for the second solution normalised by $B(0)=0$, $B(1)=1$; for the
classical $\zeta(3)$ pair we write $a_n$ and $b_n$, with the relation $b_n=6B(n)$. Digits:
for $n=ap+r$ with $0\le r<p$ we call $a=\lfloor n/p\rfloor$ the high digit and $r$ the low
digit, and similarly $k=bp+s$. Finally $\bin nk=0$ whenever $k>n$ or $k<0$, which we use
without comment.

\section{The weight-three theorem}
\label{sec:weight3}

This section is self-contained. Its only inputs are Lucas' theorem on binomial
coefficients modulo $p$ and Kummer's carry criterion for $\vp\bin NM$.

\subsection{The pair, and the harmonic decomposition}

Write
\begin{equation}\label{eq:Ank}
   A(n,k)\;=\;\bin nk^2\bin{n+k}k^2,
   \qquad
   u(n,m)\;=\;\frac{(-1)^{m-1}}{2m^3\bin nm\bin{n+m}m},
   \qquad
   c(n,k)\;=\;\HH3n+\sum_{m=1}^{k}u(n,m),
\end{equation}
so that $a_n=\sum_{k=0}^nA(n,k)$ and $b_n=\sum_{k=0}^nA(n,k)\,c(n,k)$ as in
\eqref{eq:apery-a}, \eqref{eq:apery-b}. Since $\HH3n$ does not depend on $k$, we may split
off the constant letter:
\begin{equation}\label{eq:split}
   b_n\;=\;\HH3n\,a_n\;+\;W(n),
   \qquad
   W(n)\;:=\;\sum_{k=0}^{n}A(n,k)\sum_{m=1}^{k}u(n,m).
\end{equation}
Both terms of \eqref{eq:split} will be treated separately, and the theorem is the sum of
the two congruences \eqref{eq:Hpart} and \eqref{eq:Wpart} below.

We shall also need the tail sums of the summand,
\begin{equation}\label{eq:Tdef}
  T(n,m)\;:=\;\sum_{k=m}^{n}A(n,k)\;\in\;\Z\qquad(0\le m\le n+1),
\end{equation}
with the convention $T(n,m)=0$ for $m>n$.

\begin{lemma}[interchange]\label{lem:interchange}
For every $n\ge0$,
\[
  W(n)\;=\;\sum_{m=1}^{n}u(n,m)\,T(n,m).
\]
\end{lemma}

\begin{proof}
All sums are finite, so we may reorder freely. The $k=0$ term of $W(n)$ has an empty inner
sum and contributes $0$; the remaining terms are indexed by the triangular region
$\{(m,k):1\le m\le k\le n\}$, over which
\[
  W(n)=\sum_{k=1}^{n}\sum_{m=1}^{k}A(n,k)u(n,m)
       =\sum_{\substack{1\le m\le k\le n}}A(n,k)u(n,m)
       =\sum_{m=1}^{n}u(n,m)\sum_{k=m}^{n}A(n,k),
\]
which is the assertion. Note that every $u(n,m)$ occurring has $1\le m\le n$, so
$\bin nm\ne0$ and $\bin{n+m}m\ne0$ and no division by zero occurs.
\end{proof}

\begin{remark}
Lemma~\ref{lem:interchange} is one of the three steps flagged as compressed in the
referee pass of the first write-up; it is elementary, but it is the definition of the
object $T(n,m)$ on which the whole $W$-part rests, so we have stated it as a lemma.
\end{remark}

\subsection{Lucas and Kummer inputs}

\begin{lemma}[Lucas step]\label{lem:lucas-step}
Let $p$ be prime, $a,b\ge0$ and $0\le r,s<p$. Then
$\bin{ap+r}{bp+s}\equiv\bin ab\bin rs\pmod p$.
\end{lemma}

\begin{proof}
The base-$p$ expansion of $ap+r$ has last digit $r$ and higher digits those of $a$; that
of $bp+s$ has last digit $s$ and higher digits those of $b$. Lucas' theorem gives
$\bin{ap+r}{bp+s}\equiv\bin rs\prod_{j\ge1}\bin{a_{j-1}}{b_{j-1}}$ where $a_j$, $b_j$ are
the digits of $a$, $b$; and $\prod_{j\ge1}\bin{a_{j-1}}{b_{j-1}}\equiv\bin ab$ by Lucas'
theorem again.
\end{proof}

\begin{lemma}[carry annihilation]\label{lem:carry-ann}
Fix $a,b\ge0$ and $0\le r,s<p$, and put $n=ap+r$, $k=bp+s$. Then
\[
  \bin{n+k}{k}\;\equiv\;
  \begin{cases}
    \bin{a+b}{b}\bin{r+s}{s} & \text{if }r+s<p,\\[2pt]
    0 & \text{if }r+s\ge p,
  \end{cases}
  \pmod p .
\]
\end{lemma}

\begin{proof}
We have $n+k=(a+b)p+(r+s)$. If $r+s<p$ this is a legitimate digit decomposition and
Lemma~\ref{lem:lucas-step} applies directly. If $r+s\ge p$, write
$n+k=(a+b+1)p+(r+s-p)$ with $0\le r+s-p<p$; Lemma~\ref{lem:lucas-step} then gives
$\bin{n+k}{k}\equiv\bin{a+b+1}{b}\bin{r+s-p}{s}$, and $r<p$ forces $r+s-p<s$, so
$\bin{r+s-p}{s}=0$.
\end{proof}

\begin{lemma}[summand factorisation]\label{lem:summand}
With $n=ap+r$, $k=bp+s$ as above,
\[
   A(n,k)\;\equiv\;[\,r+s<p\,]\cdot\bin ab^2\bin{a+b}b^2\cdot\bin rs^2\bin{r+s}s^2
   \pmod p,
\]
where $[\,\cdot\,]$ is the indicator of a condition. In particular $A(n,k)\equiv0\pmod p$
whenever $s>r$ or $r+s\ge p$.
\end{lemma}

\begin{proof}
$A(n,k)=\bin nk^2\bin{n+k}k^2$. By Lemma~\ref{lem:lucas-step},
$\bin nk^2\equiv\bin ab^2\bin rs^2$; by Lemma~\ref{lem:carry-ann},
$\bin{n+k}k^2\equiv[\,r+s<p\,]\bin{a+b}b^2\bin{r+s}s^2$. Multiply. The last sentence
follows since $\bin rs=0$ for $s>r$.
\end{proof}

The next two lemmas are the entire Kummer ledger, and they are where the hypothesis
$a<p$, i.e. $n<p^2$, is used.

\begin{lemma}[borrow count]\label{lem:borrow}
Let $p$ be prime, $n=ap+r$ with $0\le a,r<p$, and let $0\le m\le n$ have digits
$m=m_1p+m_0$, $0\le m_0<p$. Then
\[
  \vp\bin nm\;=\;[\,m_0>r\,]\;\in\;\{0,1\}.
\]
\end{lemma}

\begin{proof}
By Kummer's theorem, $\vp\bin nm$ is the number of borrows in the base-$p$ subtraction
$n-m$. Since $n<p^2$ and $m\le n$, both $n$ and $m$ have at most two base-$p$ digits, so
there are at most two borrow positions, and $m_1\le a$.

At position $0$ a borrow occurs precisely when $m_0>r$. At position $1$ a borrow occurs
precisely when $m_1+[\,m_0>r\,]>a$. We check that this never happens.

\emph{Case $m_0\le r$.} Then position $1$ requires $m_1>a$, contradicting $m_1\le a$.

\emph{Case $m_0>r$.} We claim $m_1\le a-1$. Indeed, if $m_1\ge a$ then
$m=m_1p+m_0\ge ap+m_0>ap+r=n$, contradicting $m\le n$. Hence $m_1+1\le a$ and there is no
borrow at position $1$.

There is no position $\ge2$, so the borrow count is $[\,m_0>r\,]$.
\end{proof}

\begin{remark}
The step ``$m_0>r\Rightarrow m_1\le a-1$'' is the second of the three compressions flagged
by the referee pass; in the first write-up it was abbreviated to ``the second borrow would
need $m_1=a$ with $m_0>r$, i.e.\ $m>n$''.
\end{remark}

\begin{lemma}[carry count]\label{lem:carrycount}
With $n=ap+r$, $0\le a,r<p$, and $0\le m=m_1p+m_0\le n$, put
\[
  c_0=[\,r+m_0\ge p\,],\qquad c_1=[\,a+m_1+c_0\ge p\,].
\]
Then $\vp\bin{n+m}{m}=c_0+c_1\le 2$.
\end{lemma}

\begin{proof}
By Kummer's theorem this valuation is the number of carries in the base-$p$ addition
$n+m$. The digits of $n$ are $(a,r)$ and those of $m$ are $(m_1,m_0)$, all in $[0,p)$. A
carry out of position $0$ occurs iff $r+m_0\ge p$, i.e.\ iff $c_0=1$; a carry out of
position $1$ occurs iff $a+m_1+c_0\ge p$, i.e.\ iff $c_1=1$. At position $2$ the digits
are $0$ and $0$ and the incoming carry is $c_1\le1<p$, so no further carry is produced.
\end{proof}

\begin{lemma}[level-$a$ carry]\label{lem:levelcarry}
If $0\le b\le a<p$ and $a+b\ge p$ then $\vp\bin{a+b}b\ge1$, hence $\vp A(a,b)\ge2$.
Consequently, if $0\le j\le a<p$ and $a+j\ge p$ then
\[
   \vp\,T(a,j)\;\ge\;2 .
\]
\end{lemma}

\begin{proof}
$a,b<p$ and $a+b\ge p$ means the base-$p$ addition of $a$ and $b$ carries out of position
$0$, so $\vp\bin{a+b}b\ge1$ by Kummer, and $A(a,b)=\bin ab^2\bin{a+b}b^2$ has valuation at
least $2$. In $T(a,j)=\sum_{b=j}^{a}A(a,b)$ every index satisfies $b\ge j$, hence
$a+b\ge a+j\ge p$, so every summand has valuation $\ge2$.
\end{proof}

\subsection{The fibre lemma, and the first row}

The following lemma does the combinatorial work three times over: it yields the first-row
Lucas congruence, and it yields the two ``tail-factorisation'' facts needed for the second
row. Its proof contains the product-completion argument in full; this is the third of the
compressions flagged in the referee pass, and it is the one that genuinely needs care,
because the true summation region is \emph{not} a product set.

\begin{lemma}[fibre lemma]\label{lem:fibre}
Let $p$ be prime, $n=ap+r$ with $a\ge0$ and $0\le r<p$. For $0\le\beta\le a+1$ set
\[
   T_\beta(n)\;:=\;\sum_{\substack{0\le k\le n\\ \lfloor k/p\rfloor\ge\beta}}A(n,k).
\]
Then
\[
   T_\beta(n)\;\equiv\;T(a,\beta)\cdot a_r\pmod p ,
   \qquad T(a,\beta)=\sum_{b=\beta}^{a}A(a,b).
\]
\end{lemma}

\begin{proof}
Write each $k$ in the range uniquely as $k=bp+s$ with $0\le s<p$; then $b=\lfloor
k/p\rfloor$ and $0\le k\le n=ap+r<(a+1)p$ forces $\beta\le b\le a$. By
Lemma~\ref{lem:summand},
\begin{equation}\label{eq:fibre1}
   T_\beta(n)\;\equiv\;
   \sum_{b=\beta}^{a}\ \sum_{\substack{0\le s<p\\ r+s<p\\ bp+s\le n}}
   \bin ab^2\bin{a+b}b^2\,\bin rs^2\bin{r+s}s^2 \pmod p .
\end{equation}
The constraint $bp+s\le n=ap+r$ is vacuous for $b<a$, and for $b=a$ it reads $s\le r$.
Now, all summands with $b=a$ and $s>r$ carry the factor $\bin rs^2=0$ and therefore vanish
\emph{identically}, not merely modulo $p$. Hence they may be adjoined to the sum
\eqref{eq:fibre1} without changing it, and after this adjunction the index region is the
product set
\[
  \{\beta\le b\le a\}\times\{0\le s<p:\;r+s<p\},
\]
in which the two constraints involve $b$ and $s$ separately. The double sum therefore
factors:
\begin{equation}\label{eq:fibre2}
  T_\beta(n)\;\equiv\;
  \Bigl(\sum_{b=\beta}^{a}\bin ab^2\bin{a+b}b^2\Bigr)
  \Bigl(\sum_{\substack{0\le s<p\\ r+s<p}}\bin rs^2\bin{r+s}s^2\Bigr)\pmod p .
\end{equation}
The first factor is $T(a,\beta)$ --- an equality of integers, by the definition
\eqref{eq:Tdef} of $T$ applied at level $a$.

For the second factor we claim
\begin{equation}\label{eq:fibre3}
   \sum_{\substack{0\le s<p\\ r+s<p}}\bin rs^2\bin{r+s}s^2\;\equiv\;a_r \pmod p .
\end{equation}
Indeed $a_r=\sum_{s=0}^{r}\bin rs^2\bin{r+s}s^2$. Terms with $s>r$ have $\bin rs=0$, so
the range may be taken to be $0\le s<p$. It remains to see that the terms with $s\le r$
and $r+s\ge p$ vanish modulo $p$: for those,
$\bin{r+s}s=\bin{1\cdot p+(r+s-p)}{0\cdot p+s}\equiv\bin10\bin{r+s-p}s=0$ by
Lemma~\ref{lem:lucas-step}, because $0\le r+s-p<s$. This proves \eqref{eq:fibre3} and,
with \eqref{eq:fibre2}, the lemma.
\end{proof}

\begin{theorem}[first row; Lucas congruence for $a_n$]\label{thm:arow}
\Proved\ Let $p$ be prime. For every $a\ge0$ and every $0\le r<p$,
\[
   a_{ap+r}\;\equiv\;a_a\,a_r\pmod p ,
\]
and hence $a_n\equiv\prod_i a_{n_i}\pmod p$ for the base-$p$ expansion $n=\sum_in_ip^i$.
\end{theorem}

\begin{proof}
Take $\beta=0$ in Lemma~\ref{lem:fibre}: the left-hand side is $T_0(n)=\sum_{k=0}^nA(n,k)=a_n$
and the right-hand side is $T(a,0)a_r=a_a\,a_r$. The product form follows by induction on
the number of base-$p$ digits of $a$, the two-digit statement being available for every
$a\ge0$.
\end{proof}

Theorem~\ref{thm:arow} is classical \cite{Gessel1982,MalikStraub,Straub2301};
we have reproved it because the fibre lemma it rests on is used twice more below, and
because the whole point of this paper is what happens one row down.

\begin{corollary}[tail factorisation]\label{cor:tailfact}
Let $n=ap+r$ with $a\ge0$, $0\le r<p$, and let $1\le m\le n$ have digits $m=m_1p+m_0$.
\begin{enumerate}[label=\textup{(\roman*)}]
\item If $m_0>r$ then $T(n,m)\equiv T(a,m_1+1)\,a_r\pmod p$.
\item If $m_0=0$, say $m=jp$ with $1\le j\le a$, then $T(n,jp)\equiv T(a,j)\,a_r\pmod p$.
\end{enumerate}
\end{corollary}

\begin{proof}
(ii) The condition $k\ge jp$ is exactly $\lfloor k/p\rfloor\ge j$, so $T(n,jp)=T_j(n)$ and
Lemma~\ref{lem:fibre} applies verbatim.

(i) Here $T(n,m)$ and $T_{m_1+1}(n)$ differ only by the terms with $m\le k<(m_1+1)p$,
i.e.\ by those $k$ with $\lfloor k/p\rfloor=m_1$ and low digit $s=k-m_1p\ge m_0>r$. By the
last sentence of Lemma~\ref{lem:summand}, each such term is $\equiv0\pmod p$. Hence
$T(n,m)\equiv T_{m_1+1}(n)\equiv T(a,m_1+1)a_r\pmod p$.
\end{proof}

\subsection{The constant letter}

\begin{lemma}[$H$-part]\label{lem:Hpart}
Let $p\ge3$, $n=ap+r$ with $0\le a,r<p$. Then
\[
   p^3\HH3n\;\equiv\;\HH3a \pmod{p^3},
\]
and consequently
\begin{equation}\label{eq:Hpart}
   p^3\,\HH3n\,a_n\;\equiv\;\HH3a\,a_a\,a_r\pmod p .
\end{equation}
\end{lemma}

\begin{proof}
The indices $m\le n$ divisible by $p$ are $m=jp$ for $1\le j\le\lfloor n/p\rfloor=a$, and
they contribute $\sum_{j=1}^{a}(jp)^{-3}=p^{-3}\HH3a$. Writing
$S=\sum_{m\le n,\;p\nmid m}m^{-3}\in\Zp$ we get $\HH3n=S+p^{-3}\HH3a$, whence
$p^3\HH3n=p^3S+\HH3a\equiv\HH3a\pmod{p^3}$. Multiplying by the integer $a_n$ preserves the
congruence modulo $p^3$, so $p^3\HH3n\,a_n\equiv\HH3a\,a_n\pmod{p^3}$, in particular
modulo $p$. Finally $\HH3a\in\Zp$ because $a<p$, and $a_n\equiv a_aa_r\pmod p$ by
Theorem~\ref{thm:arow}; multiplying the latter congruence by $\HH3a\in\Zp$ gives
\eqref{eq:Hpart}.
\end{proof}

\subsection{The singular layer}

The terms of $W(n)$ indexed by $m$ with $p\nmid m$ can carry a pole of order up to $3$,
and it is exactly the perfect-square shape of $A(n,k)$ that keeps them under control after
multiplication by $p^3$.

\begin{lemma}[singular layer]\label{lem:V}
Let $p\ge5$, $n=ap+r$ with $1\le a<p$, $0\le r<p$. For every $m$ with $1\le m\le n$ and
$p\nmid m$,
\[
   \vp\bigl(p^3\,u(n,m)\,T(n,m)\bigr)\;\ge\;1 .
\]
Consequently $\displaystyle\sum_{1\le m\le n,\;p\nmid m}p^3u(n,m)T(n,m)\equiv0\pmod p$.
\end{lemma}

\begin{proof}
Since $p\ge5$ and $p\nmid m$ we have $\vp(2m^3)=0$, so with
$CC:=\bin nm\bin{n+m}m$,
\begin{equation}\label{eq:ledger}
   \vp\bigl(p^3u(n,m)T(n,m)\bigr)\;=\;3-\vp(CC)+\vp\,T(n,m).
\end{equation}
By Lemmas~\ref{lem:borrow} and \ref{lem:carrycount},
$\vp(CC)=[\,m_0>r\,]+c_0+c_1\le3$.

If $\vp(CC)\le2$ then \eqref{eq:ledger} is $\ge3-2+0=1$, since $T(n,m)$ is an integer.

If $\vp(CC)=3$ then necessarily $[\,m_0>r\,]=c_0=c_1=1$. From $m_0>r$ and $c_1=1$, i.e.\
$a+m_1+1\ge p$, we get: Corollary~\ref{cor:tailfact}(i) applies, so
$T(n,m)\equiv T(a,m_1+1)a_r\pmod p$; and $m_1+1\ge p-a$, so every index $b\ge m_1+1$
occurring in $T(a,m_1+1)$ satisfies $a+b\ge a+m_1+1\ge p$, whence $\vp A(a,b)\ge2$ by
Lemma~\ref{lem:levelcarry} and $T(a,m_1+1)\equiv0\pmod p$. Therefore
$\vp T(n,m)\ge1$ and \eqref{eq:ledger} is $\ge3-3+1=1$.
\end{proof}

\begin{remark}\label{rem:square}
Lemma~\ref{lem:V} is the load-bearing step, and the only place where the \emph{square}
shape of $A(n,k)=\bin nk^2\bin{n+k}k^2$ is essential: a single Kummer carry buys
$\vp\ge2$, which is precisely the slack that rescues the case $\vp(CC)=3$. A summand with
an unsquared wide binomial would buy only $\vp\ge1$ and the ledger would not close. This
is a concrete obstruction for any port of the argument to a non-square summand.
\end{remark}

\subsection{The \texorpdfstring{$p$}{p}-divisible layer}

\begin{lemma}[$W$-part]\label{lem:Wpart}
Let $p\ge5$, $n=ap+r$ with $1\le a<p$, $0\le r<p$. Then
\begin{equation}\label{eq:Wpart}
   p^3\,W(n)\;\equiv\;a_r\,W(a)\pmod p .
\end{equation}
\end{lemma}

\begin{proof}
By Lemma~\ref{lem:interchange} and Lemma~\ref{lem:V},
\[
   p^3W(n)\;\equiv\;\sum_{j=1}^{a}\tau_j^{(n)}\pmod p ,
   \qquad
   \tau^{(n)}_j:=p^3u(n,jp)\,T(n,jp)
   =\frac{(-1)^{j-1}\,T(n,jp)}{2j^3\bin n{jp}\bin{n+jp}{jp}},
\]
where we used $p^3/(jp)^3=j^{-3}$ and $(-1)^{jp-1}=(-1)^{j-1}$ ($p$ odd). Likewise, by
Lemma~\ref{lem:interchange} at level $a$,
\[
   W(a)=\sum_{j=1}^{a}\tau_j,
   \qquad
   \tau_j:=u(a,j)T(a,j)=\frac{(-1)^{j-1}T(a,j)}{2j^3\bin aj\bin{a+j}j} .
\]
We compare $\tau^{(n)}_j$ with $a_r\tau_j$ for each $j$ separately. Throughout,
Lemma~\ref{lem:lucas-step} gives
\begin{equation}\label{eq:lucas-jp}
   \bin n{jp}=\bin{ap+r}{jp+0}\equiv\bin aj\bin r0=\bin aj,
   \quad
   \bin{n+jp}{jp}=\bin{(a+j)p+r}{jp+0}\equiv\bin{a+j}j \pmod p,
\end{equation}
the second because $n+jp=(a+j)p+r$ with $0\le r<p$ is a legitimate digit decomposition,
Lemma~\ref{lem:lucas-step} being valid for arbitrary high digits. Note also
$\vp\bin n{jp}=0$ by Lemma~\ref{lem:borrow} (here $m_0=0\le r$) and
$\vp\bin{n+jp}{jp}=c_0+c_1=0+[\,a+j\ge p\,]$ by Lemma~\ref{lem:carrycount}, while
$\vp\bin aj=0$ and $\vp\bin{a+j}j=[\,a+j\ge p\,]$ since $j\le a<p$.

\emph{Case $a+j<p$.} Then all four binomials in sight are $p$-units, and by
\eqref{eq:lucas-jp} the two at level $n$ are congruent modulo $p$ to the two at level
$a$; a $p$-unit congruence may be inverted modulo $p$. Combining with
Corollary~\ref{cor:tailfact}(ii), $T(n,jp)\equiv T(a,j)a_r$, we obtain
\[
   \tau^{(n)}_j\;\equiv\;\frac{(-1)^{j-1}T(a,j)\,a_r}{2j^3\bin aj\bin{a+j}j}
   \;=\;a_r\,\tau_j\pmod p .
\]

\emph{Case $a+j\ge p$.} We show that both $\tau^{(n)}_j$ and $\tau_j$ are $\equiv0\pmod
p$. For $\tau_j$: $\vp\tau_j=\vp T(a,j)-\vp\bin aj-\vp\bin{a+j}j\ge2-0-1=1$ by
Lemma~\ref{lem:levelcarry}. For $\tau^{(n)}_j$: we have
$\vp\tau^{(n)}_j=\vp T(n,jp)-0-1$, and we claim $\vp T(n,jp)\ge2$. Indeed every $k$ in the
range of $T(n,jp)=\sum_{k\ge jp}A(n,k)$ has $b:=\lfloor k/p\rfloor\ge j$, so with
$s=k-bp$ the carry count of Lemma~\ref{lem:carrycount} for $\bin{n+k}k$ satisfies
$c_1=[\,a+b+c_0\ge p\,]\ge[\,a+j\ge p\,]=1$, whence $\vp\bin{n+k}k\ge1$ and
$\vp A(n,k)\ge2$. Therefore $\vp\tau_j^{(n)}\ge2-1=1$.

Summing over $j$ and discarding the terms with $a+j\ge p$ from both sides,
\[
   p^3W(n)\equiv\sum_{j=1}^{a}\tau^{(n)}_j
   \equiv\sum_{\substack{1\le j\le a\\ a+j<p}}a_r\tau_j
   \equiv a_r\sum_{j=1}^{a}\tau_j=a_rW(a)\pmod p . \qedhere
\]
\end{proof}

\subsection{The theorem}

%% FIXED [MINOR x3]: (i) the range was $0\le a<p$ but Lemmas \ref{lem:V} and \ref{lem:Wpart},
%% FIXED which the proof invokes, both require $1\le a<p$; the $a=0$ case is true but unused, and
%% FIXED it is now excluded. (ii) the identity $W(a)=\sum_j u(a,j)T(a,j)$ used in the first
%% FIXED paragraph is Lemma \ref{lem:interchange} and was uncited. (iii) the second paragraph
%% FIXED appealed to the internals of the proof of Lemma \ref{lem:Wpart}; the statement it needs
%% FIXED is now made explicitly.
\begin{lemma}[local integrality]\label{lem:integrality}
Let $p\ge5$ and $1\le a<p$, $0\le r<p$, $n=ap+r$. Then $b_a\in\Zp$ and $p^3b_n\in\Zp$.
\end{lemma}

\begin{proof}
For $b_a=\HH3a\,a_a+W(a)$: the letter $\HH3a$ is $p$-integral because $a<p$, and by
Lemma~\ref{lem:interchange} $W(a)=\sum_{j=1}^{a}u(a,j)T(a,j)$, whose terms have
$\vp\bigl(u(a,j)T(a,j)\bigr)=-[\,a+j\ge p\,]+\vp T(a,j)$, which is $\ge0$ when $a+j<p$ and
$\ge-1+2=1$ when $a+j\ge p$ by Lemma~\ref{lem:levelcarry}.

For $p^3b_n=p^3\HH3n\,a_n+p^3W(n)$: the first term is $p$-integral by
Lemma~\ref{lem:Hpart}. For the second, use Lemma~\ref{lem:interchange} again and split the
range of $m$. The terms with $p\nmid m$ have $\vp\bigl(p^3u(n,m)T(n,m)\bigr)\ge1$ by
Lemma~\ref{lem:V}. The terms with $m=jp$ have $\vp\bigl(p^3u(n,jp)T(n,jp)\bigr)\ge0$: this is
the content of the two cases $a+j<p$ and $a+j\ge p$ established inside the proof of
Lemma~\ref{lem:Wpart}, where the first gives $\vp\ge0$ directly and the second gains the two
powers of Lemma~\ref{lem:levelcarry} against a pole of order at most $1$.
\end{proof}

\begin{theorem}[second row, weight three]\label{thm:brow}
\Proved\ Let $p\ge5$ be prime, $1\le a<p$, $0\le r<p$ and $n=ap+r$. Then both sides below
lie in $\Zp$ and
\[
   p^3\,b_{ap+r}\;\equiv\;b_a\,a_r\pmod p .
\]
\end{theorem}

\begin{proof}
Local integrality is Lemma~\ref{lem:integrality}. By \eqref{eq:split},
$p^3b_n=p^3\HH3n\,a_n+p^3W(n)$, and by \eqref{eq:Hpart} and \eqref{eq:Wpart},
\[
  p^3b_n\;\equiv\;\HH3a\,a_a\,a_r+a_r\,W(a)
  \;=\;a_r\bigl(\HH3a\,a_a+W(a)\bigr)\;=\;a_r\,b_a \pmod p . \qedhere
\]
\end{proof}

\begin{remark}[$p=5$ is not exceptional]\label{rem:p5}
The proof uses only that $p$ is odd, that $p\nmid2$, and that $1\le j\le a<p$ implies
$p\nmid j$; it therefore applies verbatim at $p=5$. This matters, because at $p=5$ one has
$a_1=5\equiv0$ and $a_3=1445\equiv0\pmod5$, so the \emph{ratio} form
$p^3(b_n/a_n)\equiv b_a/a_a$ acquires a pole (of order as low as $-1$) at the cells where
a base-$p$ digit of $a$ meets a zero of $a_\bullet$. Theorem~\ref{thm:brow} never divides
by $a_a$, so it sees no pole; the pole is an artefact of dividing, not a defect of the
congruence. Zeros of $a_\bullet$ modulo $p$ occur for $r\in\{1,3\}$ ($p=5$), $\{5\}$
($p=11$), $\{3,13\}$ ($p=17$), $\{8,10\}$ ($p=19$), $\{8,22\}$ ($p=31$), and for none of
$r$ when $p=7,13,23,29$ \VerifiedR{$5\le p\le31$}.
\end{remark}

\begin{remark}[independent verification]\label{rem:verify3}
Theorem~\ref{thm:arow} and Theorem~\ref{thm:brow}, together with every lemma above and
every case split in their proofs, were checked by an independent exact-arithmetic
implementation --- one not derived from the write-up --- with zero failures: the theorems
for $p\in\{5,7,11,13,17\}$ and all cells $(a,r)$; the fibre and tail factorisations for
$p\in\{5,7,11\}$ and all $(a,r,m)$; the Kummer counts of Lemmas~\ref{lem:borrow} and
\ref{lem:carrycount} for $p\in\{5,7,11,13\}$ and all $p\nmid m$; and the ledger of
Lemma~\ref{lem:V}, whose minimum over all $p\nmid m$ is exactly $1$, so the bound is
sharp. The dangerous branch $\vp(CC)=3$ was confirmed on all $14\,292$ of its occurrences
for $p\le19$, and $\vp(CC)$ was never observed to exceed $3$.
\end{remark}

\begin{remark}[what is \emph{not} proved here]\label{rem:not-proved}
Theorem~\ref{thm:brow} is a single-digit, modulo-$p$ statement. The multi-digit master
form and the sharpening from modulo $p$ to modulo $p^3$ are \Verified\ only; see
Section~\ref{sec:depth}. The single-digit case is the base case of a digit-by-digit
induction, but the induction step is not available, because the natural inductive
statement is the ratio form and the ratio form has poles (Remark~\ref{rem:p5}).
\end{remark}

\section{The master depth law}
\label{sec:depth}

Theorem~\ref{thm:brow} is a modulo-$p$ statement about a single base-$p$ digit. The
numerical evidence says that the true statement is stronger in three independent
directions at once: it holds modulo $p^3$, it holds for all $n$ (not only $n<p^2$), and
it needs no $p$-unit hypothesis anywhere. We state it as a conjecture and we do not
upgrade it.

\subsection{The weight-three master form}

For $n\ge1$ put $q=\lfloor n/p\rfloor$ and
\begin{equation}\label{eq:Edef}
   E(n):=\vp\bigl(p^3\,b_n\,a_q-b_q\,a_n\bigr).
\end{equation}
This quantity is defined for every $n$ and every $p$ --- no division occurs, so the
$p=5$ and exceptional-digit cells are not special.

\begin{conjecture}[master depth law, weight three]\label{conj:master3}
For every prime $p\ge5$ and every $n\ge1$, with $q=\lfloor n/p\rfloor$,
\[
   p^3\,b_n\,a_q\;\equiv\;b_q\,a_n \pmod{p^3},
   \qquad\text{i.e.}\qquad E(n)\ge3 .
\]
\end{conjecture}

\noindent\textbf{Evidence.} \VerifiedR{zero failures; all primes $5\le p\le31$; all
$n\in[1,320]$; exact rational arithmetic}. This range covers the single-digit cells
$n<p^2$, the multi-digit range, and $p=5$. The floor is attained: $\min_nE(n)=3$ exactly,
for every one of the nine primes. On $p$-unit cells the statement is equivalent to the
ratio descent
\[
   p^3\,\frac{b_n}{a_n}\;\equiv\;\frac{b_q}{a_q}\pmod{p^3},
\]
which iterates down the whole base-$p$ expansion; and the minimum of
$\vp\bigl(p^3b_n/a_n-b_q/a_q\bigr)$ over $p$-unit single-digit cells is likewise exactly
$3$, for every prime $7\le p\le31$, with no cell of depth $\le2$
\VerifiedR{$7\le p\le31$}.

Restricted to single-digit cells $n=ap+r$ ($1\le a<p$, $0\le r<p$) the quantity
\eqref{eq:Edef} is $E(a,r)=\vp(p^3b_na_a-b_aa_n)$, and its distribution is remarkably
rigid: overwhelmingly $3$, with a handful of sporadic boosts.

\begin{table}[ht]
\centering
\caption{Distribution of the single-digit depth
$E(a,r)=\vp\bigl(p^3b_{ap+r}a_a-b_aa_{ap+r}\bigr)$ over all cells $1\le a<p$,
$0\le r<p$. \VerifiedR{exact; zero cells with $E<3$}.}
\label{tab:Etally}
\begin{tabular}{@{}rcl@{}}
\toprule
$p$ & $\min E$ & $E$-value tally over single-digit cells ($E$:\ count)\\
\midrule
 5 & 3 & $3{:}18,\ 4{:}1,\ 6{:}1$\\
 7 & 3 & $3{:}35,\ 4{:}5,\ 5{:}1,\ 6{:}1$\\
11 & 3 & $3{:}104,\ 4{:}4,\ 6{:}2$\\
13 & 3 & $3{:}147,\ 4{:}7,\ 5{:}1,\ 6{:}1$\\
17 & 3 & $3{:}260,\ 4{:}10,\ 5{:}1,\ 6{:}1$\\
19 & 3 & $3{:}302,\ 4{:}35,\ 5{:}4,\ 6{:}1$\\
23 & 3 & $3{:}478,\ 4{:}25,\ 5{:}2,\ 6{:}1$\\
29 & 3 & $3{:}774,\ 4{:}35,\ 5{:}2,\ 6{:}1$\\
31 & 3 & $3{:}912,\ 4{:}16,\ 5{:}1,\ 6{:}1$\\
\bottomrule
\end{tabular}
\end{table}

\subsection{Fine structure of the boosts}

%% FIXED [MINOR]: the reflection (a,r) -> (p-1-a, p-1-r) sends a = p-1 to a = 0, which is
%% FIXED outside the cell range 1 <= a < p, so it cannot pair the edge and corner boosts.
The cells with $E>3$ are not noise, but they do not follow a clean law either. They
concentrate at extreme digits, and the interior ones are symmetric under the reflection
$(a,r)\mapsto(p-1-a,\,p-1-r)$ --- a symmetry which says nothing about the $a=p-1$ boosts,
since it carries them to $a=0$, outside the range of the table \VerifiedR{$5\le p\le31$}:
\begin{itemize}[leftmargin=1.4em]
%% FIXED [MAJOR-3]: "E = 5 at the edge" was false at p=5 (E=4) and p=11 (E=6), contradicting
%% FIXED Table \ref{tab:Etally} above, which lists no E=5 cell at either prime, and the bullet
%% FIXED immediately above it. Recomputed exactly here for 5<=p<=31: E(edge) = 4, 5, 6, 5, 5, 5,
%% FIXED 5, 5, 5 at p = 5, 7, 11, 13, 17, 19, 23, 29, 31. The correction was already on file at
%% FIXED work/VERIFY_ZETA3_PROOF.md:87 ("E(edge (p-1,0)) in {4,5,6}") and was dropped in favour
%% FIXED of the uncorrected WARMUP_ZETA3_DWORK.md:125.
%% FIXED Also fixed: at p=11 the corner and edge cells BOTH have E=6 (Table \ref{tab:Etally}
%% FIXED shows 6:2), so the boost is duplicated, not split.
\item $E=6$ at the corner $(a,r)=(p-1,p-1)$, i.e.\ $n=p^2-1$, at every prime in range;
\item $E\in\{4,5,6\}$ at the edge $(a,r)=(p-1,0)$, i.e.\ $n=p(p-1)$: it is $5$ at
$p=7,13,17,19,23,29,31$, but $4$ at $p=5$ and $6$ at $p=11$. At $p=11$ the edge therefore
carries a second $E=6$ cell, which is the whole of that prime's $6{:}2$ entry in
Table~\ref{tab:Etally};
\item $E=4$ on symmetric interior and edge pairs, e.g.\ for $p=13$ at
$(3,5)$, $(3,7)$, $(9,5)$, $(9,7)$, $(12,2)$, $(12,6)$, $(12,10)$.
\end{itemize}
The robust, conjecturable content is the floor $E\ge3$; the boosts we record as data.

\subsection{Two strengthenings that the data kills}

It is worth recording two natural-looking sharpenings that are \emph{false}, because both
were conjectured at some point in this programme.

\begin{observation}\label{obs:false1}
The floor is \emph{flat}: it is independent of $\kappa=\vp\bin{2n}n$. Cells with
$\kappa=0,1,2$ all floor at $E=3$ \VerifiedR{$5\le p\le31$, $n\le320$}. In particular the
bound $E\ge\vp(a_aa_n)+1$ is false: $E-\vp(a_aa_n)$ reaches $-2$ at $p=17$.
\end{observation}

\begin{observation}\label{obs:false2}
The integrality statement $\vp(b_n)\ge-3\vp(d_n)$, equivalently $d_n^3b_n\in\Z$, holds
\VerifiedR{$5\le p\le31$, $n\le320$, zero failures}, with $\vp(d_n)=\lfloor\log_pn\rfloor$.
But the corresponding \emph{equality} $\vp(b_n/a_n)=-3\vp(d_n)$ is false: the numerators
of $b_n$ carry sporadic extra $p$-factors, for instance $7\mid62531=\operatorname{num}(b_3)$.
The correct statement is an inequality.
\end{observation}

\subsection{The weight panel}
\label{ssec:panel}

Running the same experiment at lower weight explains both why the weight-three case is
clean and where the danger lives. Let $(A,B)$ be an Ap\'ery-like pair with Ap\'ery limit
of weight $w$, and consider the master quantity $\vp\bigl(p^wB_nA_q-B_qA_n\bigr)$,
$q=\lfloor n/p\rfloor$.

\begin{observation}[congruence depth $=$ weight]\label{obs:panel}
\VerifiedR{zero failures} the master floor is $\ge w$, flat, for all $p\ge5$:
\begin{itemize}[leftmargin=1.4em]
\item $w=1$: the central Delannoy numbers $1,3,13,63,\dots$ with their companion,
  $B_n/A_n\to\log2$; $n\le300$.
\item $w=2$: the Ap\'ery pair for $\zeta(2)$, $(n+1)^2u_{n+1}=(11n^2+11n+3)u_n+n^2u_{n-1}$,
  $A=1,3,19,147,1251,\dots$; $n\le300$, floor exactly $2$ for $p\ge7$.
\item $w=3$: the pair \eqref{eq:apery-a}--\eqref{eq:apery-b}, floor exactly $3$
  (Conjecture~\ref{conj:master3}).
\end{itemize}
\end{observation}

The panel also shows what the \emph{ratio} form does at low weight, and it is instructive.
At $w=2$ the ratio descent $\vp\bigl(p^2B_n/A_n-B_q/A_q\bigr)$ has a clean floor $2$ at
all residues except the reflection centre $r=(p-1)/2$, where the behaviour splits by $p$
modulo $4$: for $p\equiv1\pmod4$ it \emph{boosts} (values $5,8$ at $p=5$; $5$ at $p=13$),
while for $p\equiv3\pmod4$ it \emph{collapses to $0$} ($p=7,11$) and the naive modulo-$p$
statement fails outright. At $w=1$ the floor is barely $1$, the centre fails at $p=7$
(depth $0$) and $p=11$ (depth $-1$), and there are sporadic dips to $0$ elsewhere ($p=13$
at $r=2,10$). Exceptional digits behave exactly as textbook: the case $p=7$, $a=3$ fails
wholesale at $w=2$ because $7^2\mid A(3)=147$; at $w=1$ the failures sit at
$(p,a)=(7,3)$ with $7\mid63$, at $(13,2)$ with $13=A(2)$, and at $(3,1)$ with $3\mid3$.

\begin{observation}[the product form is the right universal statement]\label{obs:product}
Every low-weight ratio-form failure just listed is a division artefact: it occurs exactly
where $p$ divides the Lucas factor $A(a)$. The product form
\[
   p^{w}B_{ap+r}\;\equiv\;B_a\,A_r\pmod p
\]
holds at \emph{every} previously failing site \VerifiedR{$w=2$: $p=7,11$, all $a\le3$;
$w=1$: $p=7,11,13$, all $a\le3$, including the cell $13=A(2)$; $w=3$: $p=5$, $a\le3$}.
\end{observation}

Observation~\ref{obs:product} is the reason this paper states its theorems in product form
and not in ratio form, and it is what makes Theorem~\ref{thm:brow} valid at $p=5$
(Remark~\ref{rem:p5}). Historically, the programme first recorded a ``dip to depth $2$ at
the reflection centre'' at weight three; the full sweep showed that those cells are
precisely the ratio-form pole cells --- $11\mid a_5=819005$, for example --- and that on
unit cells the depth floor is a flat $3$. The correction is recorded here because it is
the same phenomenon that governs the low-weight panel.

\subsection{The general master form}

Section~\ref{sec:families} will show that at the level of all fifteen sporadic pairs one
further correction is needed, a quadratic character $\chi$. The general statement, of
which Conjecture~\ref{conj:master3} is the case $\gamma$ (with $\chi$ trivial and $w=3$),
is:

\begin{conjecture}[master form $(\mathrm{LB}_w^\chi)$]\label{conj:master}
Let $(A,B)$ be one of the fifteen sporadic Ap\'ery-like pairs, $w$ the weight of its
Ap\'ery limit and $\chi$ the quadratic character attached to it
(Table~\ref{tab:fifteen}). Then for every prime $p\ge5$ and every $n\ge1$, with
$q=\lfloor n/p\rfloor$,
\[
   \vp\bigl(p^{w}B_nA_q-\chi(p)\,B_qA_n\bigr)\;\ge\;w .
\]
\end{conjecture}

\noindent\textbf{Evidence.} \VerifiedR{zero failures} for all fifteen pairs, every prime
$5\le p\le103$, all $n\le400$; and, for $p=5,7,11$, all $n\le1200$ --- five base-$p$
digits at $p=5$. Exact rational arithmetic throughout. The floor is exactly $w$ in every
single (sequence, prime) cell: perfectly flat, with no $\kappa$-dependence, no
centre-residue exception and no exceptional-digit exception. The runs were organised as
$15\times16$ cells for the primes $5\le p\le61$ with $n\le400$, plus $15\times9$ cells for
the primes $67\le p\le103$ with $n\le400$, plus $15\times3$ cells for $p=5,7,11$ with
$n\le1200$. In the same sweeps the single-digit twisted Lucas form
$p^wB_{ap+r}\equiv\chi(p)B_aA_r\pmod p$ and the integrality bound
$\vp(B_n)\ge-w\lfloor\log_pn\rfloor$ each recorded zero failures.

The mod-$p$, single-digit shadow of Conjecture~\ref{conj:master} is a theorem for a
specified subclass; that is the content of Section~\ref{sec:general}.

\section{The general \texorpdfstring{$\chi$}{chi}-twisted descent}
\label{sec:general}

The proof of Theorem~\ref{thm:brow} in Section~\ref{sec:weight3} is a two-page Kummer
ledger. It turns out that the ledger is an artefact of Ap\'ery's own choice of weight: if
one replaces the nested alternating sum $\sum_{m\le k}u(n,m)$ by a plain harmonic monomial
--- which Section~\ref{sec:minimal} shows is possible, and in a startlingly short form ---
the whole argument collapses to three lines, and generalises at once to all weights, to
several sequences, and to a twist by a quadratic character.

\subsection{Lemma K: Frobenius descent of a harmonic letter}

\begin{lemma}[Lemma K]\label{lem:K}
\Proved\ Let $\chi$ be a completely multiplicative arithmetic function (a Dirichlet
character; $\chi\equiv1$ is allowed), let $r\ge1$ and $y\ge0$, and set
$\Kchi ry=\sum_{m=1}^{y}\chi(m)m^{-r}$. Then, as an identity in $\Q$,
\[
   p^{r}\,\Kchi ry\;=\;\chi(p)\,\Kchi r{\lfloor y/p\rfloor}\;+\;p^{r}\,T,
   \qquad
   T:=\sum_{\substack{m\le y\\ p\nmid m}}\frac{\chi(m)}{m^{r}}\;\in\;\Zp .
\]
\end{lemma}

\begin{proof}
Split the sum according to $p\mid m$. The $p$-divisible part is
\[
   \sum_{j\le\lfloor y/p\rfloor}\frac{\chi(jp)}{(jp)^{r}}
   \;=\;\chi(p)\,p^{-r}\sum_{j\le\lfloor y/p\rfloor}\frac{\chi(j)}{j^{r}}
   \;=\;\chi(p)\,p^{-r}\,\Kchi r{\lfloor y/p\rfloor},
\]
using complete multiplicativity. The remaining terms have $p$-integral denominators, so
their sum is $T\in\Zp$. Multiply by $p^r$.
\end{proof}

\emph{This single line is the entire source of the character in Conjecture~\ref{conj:master}.}
Sequences whose second solution needs $\chi$-twisted harmonic letters acquire $\chi(p)$;
sequences whose second solution is a pure $\zeta$-harmonic monomial acquire $+1$.

\begin{corollary}[Corollary K2]\label{cor:K2}
\Proved\ Let $M=\prod_{t=1}^{s}K^{(r_t)}_{\chi_t}(x_t)$ with $\sum_tr_t=w$ and each
$K^{(r_t)}_{\chi_t}(\lfloor x_t/p\rfloor)\in\Zp$. Then
\[
   p^{w}M\;\equiv\;\Bigl(\prod_t\chi_t(p)\Bigr)\prod_tK^{(r_t)}_{\chi_t}
   \bigl(\lfloor x_t/p\rfloor\bigr)\pmod p .
\]
\end{corollary}

\begin{proof}
$p^wM=\prod_t\bigl(p^{r_t}K^{(r_t)}_{\chi_t}(x_t)\bigr)
=\prod_t\bigl(\chi_t(p)K^{(r_t)}_{\chi_t}(\lfloor x_t/p\rfloor)+p^{r_t}T_t\bigr)$
by Lemma~\ref{lem:K}. Expanding, every term other than the pure product carries a factor
$p$, and every factor is $p$-integral by hypothesis.
\end{proof}

\subsection{Theorem LB}

\begin{definition}[admissible pair]\label{def:admissible}
Let $p\ge5$. An \emph{admissible pair} consists of
\begin{itemize}[leftmargin=1.4em]
\item a summand $S(n,k)=\prod_i\bin{L_i(n,k)}{M_i(n,k)}$ with $L_i,M_i$ integer linear
  forms, and $A(n)=\sum_kS(n,k)$;
\item a weight $w(n,k)=\sum_jc_j\prod_tK^{(r_{jt})}_{\chi}(x_{jt})$, a $\Q$-combination of
  monomials of total degree $\sum_tr_{jt}=w$ in letters attached to one fixed quadratic
  character $\chi$ and to the trivial character, with the $x_{jt}$ integer linear forms in
  $(n,k)$;
\item the second row $B(n)=\sum_kS(n,k)w(n,k)$.
\end{itemize}
\end{definition}

Fix $1\le a<p$ and $0\le r,s<p$, and write $n=ap+r$, $k=bp+s$. The hypotheses are:

\begin{description}[leftmargin=3.0em,style=sameline]
\item[(H1)] \emph{Lucas/carry dichotomy.} For every $k$, either $p\mid S(n,k)$, or
  $S(n,k)\equiv S(a,b)S(r,s)\pmod p$.
\item[(H2)] \emph{Product region.} The surviving set is $\{0\le b\le a\}\times\Sigma_r$
  for some $\Sigma_r\subseteq[0,p)$, and $\sum_{s\in\Sigma_r}S(r,s)\equiv A(r)\pmod p$.
\item[(H3)] \emph{Digit compatibility.} On the surviving set,
  $\lfloor x(n,k)/p\rfloor=x(a,b)$ for every argument form $x$ occurring in $w$.
\item[(H4)] \emph{Tameness.} Every argument form satisfies $0\le x(n,k)\le n$ on the
  summation range, and $c_j\in\Zp$ for all $j$.
\item[(H5)] \emph{$\chi$-homogeneity.} Every monomial of $w$ contains exactly the same
  number $e\in\{0,1\}$ of $\chi$-letters, all other letters carrying the trivial
  character.
\end{description}

\begin{theorem}[Theorem LB]\label{thm:LB}
\Proved\ Let $p\ge5$ and let $(S,w,A,B)$ be an admissible pair satisfying
\textup{(H1)--(H5)}. Then for all $1\le a<p$ and $0\le r<p$, with $n=ap+r$,
\[
   p^{w}\,B(ap+r)\;\equiv\;\chi(p)^{e}\,B(a)\,A(r)\pmod p .
\]
\end{theorem}

\begin{proof}
By (H4), every argument satisfies $x(n,k)\le n<p^2$, hence $\lfloor x/p\rfloor\le a<p$ and
every letter $K^{(r)}(\lfloor x/p\rfloor)$ is $p$-integral; likewise $x(a,b)\le a<p$, so
%% FIXED [MINOR]: the one-line reason for $p^w w(n,k) \in \Zp$ was missing.
$w(a,b)\in\Zp$. For $p^ww(n,k)$: each letter $K^{(r)}_\chi(x)$ with $x\le n<p^2$ has
$\vp\ge-r$ by Lemma~\ref{lem:K}, a monomial $\prod_tK^{(r_{jt})}_\chi$ of total degree $w$
therefore has $\vp\ge-w$, and (H4) puts the coefficients $c_j$ in $\Zp$; hence
$p^ww(n,k)\in\Zp$ for \emph{every} $k$.
Corollary~\ref{cor:K2} then gives, for every $k$,
\begin{equation}\label{eq:LB-letter}
   p^{w}w(n,k)\;\equiv\;\chi(p)^{e}\,\widetilde w(n,k)\pmod p,
   \qquad
%% FIXED [MINOR]: the letter is the chi-twisted $K^{(r)}_\chi$; the subscript was dropped
%% FIXED (log typo at LBW_GENERAL.md:271).
   \widetilde w(n,k):=\sum_jc_j\prod_tK^{(r_{jt})}_{\chi}\bigl(\lfloor x_{jt}/p\rfloor\bigr),
\end{equation}
and (H3) upgrades $\widetilde w(n,k)=w(a,b)$ on the surviving set. Now split
\[
   p^{w}B(n)\;=\;\sum_{k\,:\,p\mid S(n,k)}S(n,k)\,p^{w}w(n,k)
   \;+\;\sum_{k\text{ surviving}}S(n,k)\,p^{w}w(n,k).
\]
The first sum is $\equiv0\pmod p$, because each $p^ww(n,k)$ is $p$-integral and each
$S(n,k)$ is divisible by $p$. In the second, substitute (H1) and
\eqref{eq:LB-letter}:
\[
   p^{w}B(n)\;\equiv\;\chi(p)^{e}\sum_{b=0}^{a}\sum_{s\in\Sigma_r}S(a,b)S(r,s)w(a,b)
   \;=\;\chi(p)^{e}\Bigl(\sum_{b=0}^{a}S(a,b)w(a,b)\Bigr)\Bigl(\sum_{s\in\Sigma_r}S(r,s)\Bigr),
\]
which by (H2) is $\equiv\chi(p)^{e}B(a)A(r)\pmod p$.
\end{proof}

\subsection{Remarks on the hypotheses}

\begin{remark}[three lines instead of two pages]\label{rem:threeLines}
With $e=0$, $S(n,k)=A(n,k)$ and the weight $w(n,k)=2\HH3n-\HH3k$ of
Theorem~\ref{thm:minimal} below, Theorem~\ref{thm:LB} \emph{reproves
Theorem~\ref{thm:brow}}. Hypotheses (H1)--(H3) for this instance are exactly
Lemmas~\ref{lem:summand} and \ref{lem:fibre}; (H4) holds because the arguments are $n$ and
$k$, both $\le n$, and the coefficients are the integers $2,-1$; (H5) is vacuous. No
Lemma~\ref{lem:V}, no Corollary~\ref{cor:tailfact}, no Kummer ledger. The entire content
of the long proof is absorbed into the choice of weight.
\end{remark}

\begin{remark}[what the mechanism actually needs]\label{rem:needs}
In the language of Section~\ref{sec:weight3}: ``summand with Lucas factorisation and carry
annihilation'' is (H1)$+$(H2), which is available for all fifteen sporadic sequences
because it is Malik--Straub's theorem \cite{MalikStraub}; ``harmonic weight of depth one''
is the shape of $w$; and the per-carry valuation bound $\vp\ge2$ of
Remark~\ref{rem:square} is \emph{not} needed --- it is replaced by (H4), tameness, which
makes $p^ww$ $p$-integral outright. Where tameness fails --- because an argument form
reaches $2n$, so that $p^w w(n,k)$ acquires a pole on the vanishing layer --- one needs a
substitute valuation bound; see Section~\ref{ssec:lemmaD}.
\end{remark}

\begin{remark}[scope: single digit, modulo $p$]\label{rem:LBscope}
Theorem~\ref{thm:LB} gives the modulo-$p$, single-digit statement. The master form of
Conjecture~\ref{conj:master} --- all $n$, modulo $p^w$ --- is \Verified\ but not proved
here. The gap is exactly the one left open at weight three
(Remark~\ref{rem:not-proved}).
\end{remark}

\begin{corollary}[integrality]\label{cor:integrality}
Iterating the ratio form of Conjecture~\ref{conj:master} over the base-$p$ digits gives
$\vp(B_n)\ge-w\lfloor\log_pn\rfloor$, i.e.\ $d_n^{w}B_n\in\Zp$ --- the coefficient-level
shadow of $p$-integrality of the mirror map, see Section~\ref{sec:discussion}. Independent
of the conjecture, this bound is \VerifiedR{15 sequences $\times$ 16 primes $\times$
$n\le400$, zero failures}.
\end{corollary}

\begin{remark}[the formalised hypotheses differ, and are weaker]\label{rem:leanH}
The Lean~4 formalisation of Theorem~\ref{thm:LB} (Section~\ref{sec:lean}) replaces
(H1)$+$(H2) by the single unconditional congruence
$S(ap+r,bp+s)\equiv S(a,b)S(r,s)\pmod p$ for \emph{all} $a,b$ together with the vanishing
$S(n,k)=0$ for $k>n$. In the carrying regime both sides vanish modulo $p$, so the
dichotomy is absorbed and no product-region argument occurs anywhere in the machine proof;
the double sum factors because a \emph{full} block $\{0\le k<(a+1)p\}$ splits, and the
extension from $\{0\le k\le n\}$ to the full block is an equality of integers rather than
a congruence. It also replaces the counting form of (H5) by the equivalent but more
flexible requirement that the product of the characters of the letters of each monomial,
evaluated at $p$, be constant.
\end{remark}

\section{The fifteen sporadic families}
\label{sec:families}

\subsection{The pairs}

Zagier's search and its Almkvist--Zudilin and Cooper extensions produce fifteen sporadic
Ap\'ery-like sequences, satisfying one of
\begin{align}
  \text{(R2)}&& (n+1)^2u_{n+1}&=(an^2+an+b)u_n-cn^2u_{n-1}, \label{eq:R2}\\
  \text{(R3)}&& (n+1)^3u_{n+1}&=(2n+1)(an^2+an+b)u_n-n(cn^2+d)u_{n-1}. \label{eq:R3}
\end{align}
Zagier's six satisfy \eqref{eq:R2}; the six Almkvist--Zudilin sequences satisfy
\eqref{eq:R3} with $d=0$, and Cooper's three satisfy \eqref{eq:R3} with $d\ne0$. The
sequence tables we use are those of Malik--Straub \cite[Tables 1--2]{MalikStraub}; all
fifteen binomial sums were \VerifiedR{$n\le11$} to reproduce the recurrence solution
exactly.

The \emph{second solution} $B(n)$ is normalised by $B(0)=0$, $B(1)=1$, with the recurrence
enforced for $n\ge1$ --- the $n=0$ step is degenerate. This is exactly Ap\'ery's
construction, and the same normalisation used by Chamberland--Straub \cite{ChamStraub}.
Rescaling $B$ by a $p$-unit rational does not affect any statement below, so for $p\ge5$
the normalisation is canonical up to harmless constants.

\begin{table}[ht]
\centering\footnotesize\setlength{\tabcolsep}{4pt}
\caption{The fifteen sporadic pairs: parameters of \eqref{eq:R2}/\eqref{eq:R3}, the
integral solution $A(n)$, the Ap\'ery limit $L=\lim B(n)/A(n)$, the weight $w$ and the
character $\chi$.}
\label{tab:fifteen}
\begin{tabular}{@{}clllccl@{}}
\toprule
\# & label & $(a,b,c[,d])$ & $A(n)$ & $L$ & $w$ & $\chi$\\
\midrule
1 & $\mathbf A$ & $(7,2,-8)$ & $\sum_k\bin nk^3$ (Franel) & $\zeta(2)/4$ & 2 & $1$\\
2 & $\mathbf B$ & $(9,3,27)$ & $\sum_k(-1)^k3^{n-3k}\bin n{3k}\frac{(3k)!}{k!^3}$ & none & 2 & $\chi_{-3}$\\
3 & $\mathbf C$ & $(10,3,9)$ & $\sum_k\bin nk^2\bin{2k}k$ & $L_{-3}(2)/2$ & 2 & $\chi_{-3}$\\
4 & $\mathbf D$ & $(11,3,-1)$ & $\sum_k\bin nk^2\bin{n+k}n$ & $\zeta(2)/5$ & 2 & $1$\\
5 & $\mathbf E$ & $(12,4,32)$ & $\sum_k\bin nk\bin{2k}k\bin{2(n-k)}{n-k}$ & $L_{-4}(2)/2$ & 2 & $\chi_{-4}$\\
6 & $\mathbf F$ & $(17,6,72)$ & $\sum_k(-1)^k8^{n-k}\bin nk\sum_l\bin kl^3$ & $\tfrac58L_{-3}(2)$ & 2 & $\chi_{-3}$\\
7 & $(\alpha)$ & $(10,4,64,0)$ & $\sum_k\bin nk^2\bin{2k}k\bin{2(n-k)}{n-k}$ (Domb) & $7\zeta(3)/24$ & 3 & $1$\\
8 & $(\gamma)$ & $(17,5,1,0)$ & $\sum_k\bin nk^2\bin{n+k}n^2$ (Ap\'ery) & $\zeta(3)/6$ & 3 & $1$\\
9 & $(\delta)$ & $(7,3,81,0)$ & $\sum_k(-1)^k3^{n-3k}\bin n{3k}\bin{n+k}n\frac{(3k)!}{k!^3}$ & none & 3 & $1$\\
10 & $(\varepsilon)$ & $(12,4,16,0)$ & $\sum_k\bin nk^2\bin{2k}n^2$ & $7\zeta(3)/32$ & 3 & $1$\\
11 & $(\zeta)$ & $(9,3,-27,0)$ & $\sum_{k,l}\bin nk^2\bin nl\bin kl\bin{k+l}n$ & $L_{-3}(3)/3$ & 3 & $\chi_{-3}$\\
12 & $(\eta)$ & $(11,5,125,0)$ & $\sum_k(-1)^k\bin nk^3\bigl[\bin{4n-5k-1}{3n}+\bin{4n-5k}{3n}\bigr]$ & none & 3 & $\chi_5$\\
13 & $s_7$ & $(13,4,-27,3)$ & $\sum_k\bin nk^2\bin{n+k}k\bin{2k}n$ & $\zeta(2)/7$ & 2 & $1$\\
14 & $s_{10}$ & $(6,2,-64,4)$ & $\sum_k\bin nk^4$ & $\zeta(2)/5$ & 2 & $1$\\
15 & $s_{18}$ & $(14,6,192,-12)$ & Cooper's $s_{18}$ & $L_{-3}(2)/2$ & 2 & $\chi_{-3}$\\
\bottomrule
\end{tabular}
\end{table}

\subsection{Limits, weights and integrality}

\begin{observation}[Ap\'ery limits]\label{obs:limits}
\Verified\ The limits in Table~\ref{tab:fifteen} were identified by two-term PSLQ at $400$
digits against a $33$-constant basis --- $\zeta(2),\zeta(3),\zeta(4)$; $L_D(2)$ and
$L_D(3)$ for $D=-3,-4,-7,-8,-11,-15,-20,-24,5,8,12,13,24$; $\log^22$, $\pi^2\log2$,
$\log^32$ --- matching to between $31$ and $401$ digits (the weakest is $\mathbf F$, whose
convergence rate is $(8/9)^n$; $s_{18}$ converges at $(3/4)^n$). The known values are
recovered: $\zeta(2)/4$ for Franel and $\zeta(2)/5$ for $\sum_k\bin nk^4$
(cf.\ \cite[eq.~(7.2)]{ChamStraub}) and $\zeta(3)/6$ for Ap\'ery. The remaining eight ---
$7\zeta(3)/24$ (Domb), $7\zeta(3)/32$ ($\varepsilon$), $L_{-3}(3)/3$ ($\zeta$),
$\zeta(2)/7$ ($s_7$), $\tfrac58L_{-3}(2)$ ($\mathbf F$), $L_{-3}(2)/2$ ($\mathbf C$ and
$s_{18}$) and $L_{-4}(2)/2$ ($\mathbf E$, half Catalan's constant) --- are computed here.
They are of the shape one expects (rational multiples of $\zeta(2)$, $L_{-3}(2)$,
$L_{-4}(2)$), and some may well appear in existing numerical tables; we make no
attribution claim.
\end{observation}

\begin{observation}[weights]\label{obs:weights}
\Verified\ Let $w$ be the minimal exponent with $d_n^wB(n)\in\Z$, computed over $n\le600$.
Then $w=2$ for all six of Zagier's \eqref{eq:R2} sequences \emph{and} for all three of
Cooper's \eqref{eq:R3} sequences with $d\ne0$, and $w=3$ for all six Almkvist--Zudilin
sequences ($d=0$). It is sharp in every case: $d_n^{w-1}B_n$ is non-integral for at least
$596$ of the $600$ values of $n$. In every case $w$ equals the motivic weight of the
Ap\'ery limit. The local form $\vp(B_n)\ge-w\lfloor\log_pn\rfloor$ holds with zero
failures over 15 sequences $\times$ 16 primes $\times$ $n\le400$.
\end{observation}

\begin{observation}[Cooper's three are lighter than their recurrence]\label{obs:cooper}
$s_7$, $s_{10}$, $s_{18}$ have weight $2$, not $3$, despite satisfying the cubic
recurrence \eqref{eq:R3}: the weight tracks the Cooper deformation $d\ne0$, not the shape
of the recurrence.
\end{observation}

\subsection{Where the untwisted law dies, and the repair}
\label{ssec:twist}

\begin{observation}[failure of the naive law]\label{obs:naive-fail}
\Verified\ Testing $\vp\bigl(p^wB_nA_q-B_qA_n\bigr)\ge w$, $q=\lfloor n/p\rfloor$, over
primes $5\le p\le43$ and $n\le300$ in exact arithmetic:
\begin{center}\small
\begin{tabular}{@{}ll@{}}
\toprule
verdict & sequences\\
\midrule
holds flat, every prime tested & $\mathbf A$, $\mathbf D$, $\alpha$, $\gamma$, $\delta$, $\varepsilon$, $s_7$, $s_{10}$ \quad(8)\\
fails at $p\equiv2\pmod3$ & $\mathbf B$, $\mathbf C$, $\mathbf F$, $\zeta$, $s_{18}$ \quad(5)\\
fails at $p\equiv3\pmod4$ & $\mathbf E$ \quad(1)\\
fails at $p\equiv\pm2\pmod5$ & $\eta$ \quad(1)\\
\bottomrule
\end{tabular}
\end{center}
The failures are total, not marginal: on the failing primes the single-digit floor is $0$
--- the difference is a $p$-unit --- and in the multi-digit range the valuation drops to
$-w(\#\text{digits}-1)$, i.e.\ as low as the integrality bound permits.
\end{observation}

The failing sets are not explained by exceptional digits ($p\mid A(a)$), nor by the centre
residue, nor by $\kappa=\vp\bin{2n}n$: they are exactly the sets of primes at which a
quadratic character takes the value $-1$. Three cross-checks isolate the character
\Verified: the passing primes of the $\chi_{-3}$ group are mixed modulo $4$ and modulo
$5$; those of $\mathbf E$ are mixed modulo $3$ and modulo $5$; and those of $\eta$, namely
$11,19,29,31,41$, are mixed modulo $3$ and modulo $4$ and are precisely the primes
$p\equiv\pm1\pmod 5$.

\begin{observation}[the repair is one sign]\label{obs:repair}
\VerifiedR{zero failures} At every cell with $\chi(p)=-1$, for all seven twisted
sequences, all primes $5\le p\le31$ with $\chi(p)=-1$ and all $n<p^3$ (up to $1500$),
\[
   \vp\bigl(p^{w}B_nA_q+B_qA_n\bigr)\;=\;w \quad\text{exactly.}
\]
\end{observation}

Observations~\ref{obs:naive-fail} and \ref{obs:repair} together are
Conjecture~\ref{conj:master}. The character is the character of the Ap\'ery limit --- a
correlation that is $12$ out of $12$ on the sequences whose limit exists --- and
Lemma~\ref{lem:K} says why: the $\chi(p)$ in the congruence is the $\chi(p)$ that comes
out of the $p$-divisible layer of a $\chi$-twisted harmonic letter.

\begin{observation}[ramified primes are benign]\label{obs:ramified}
For $\eta$ the prime $p=5$ is ramified, $\chi_5(5)=0$, and
Conjecture~\ref{conj:master} degenerates to $p^3B_nA_q\equiv0\pmod{p^3}$ --- which is
exactly the assertion that $B_n$ is $5$-integral. \VerifiedR{$n\le500$} it is:
$\min_n v_5(B_n)=0$, so the pole simply does not develop at the ramified prime, and the
floor is exactly $3$. This is the only ramified prime $\ge5$ among the fifteen; the
$\chi_{-3}$ and $\chi_{-4}$ sequences ramify only at $3$ and $2$.
\end{observation}

\subsection{Harmonic decompositions}
\label{ssec:decomp}

To apply Theorem~\ref{thm:LB} one needs a harmonic-monomial weight. We found these by
exact linear fitting: the ansatz is $B_n=\sum_kS(n,k)w(n,k)$ with $S$ the sequence's own
binomial summand and $w$ a $\Q$-combination of weight-$w$ monomials in letters $\HH rx$
(and, when needed, $K^{(r)}_{\chi,x}$) at arguments $x$ that are linear forms in $(n,k)$;
the systems were solved exactly modulo the prime $2^{61}-1$, the coefficients rationally
reconstructed, and then \emph{validated exactly over $\Q$} on held-out values of $n$.

\begin{table}[ht]
\centering\footnotesize
\caption{Harmonic-monomial weights found (all depth one, all validated exactly over $\Q$
on held-out $n$). $S$ is the summand from Table~\ref{tab:fifteen}.}
\label{tab:decomp}
\begin{tabular}{@{}llll@{}}
\toprule
seq. & $w$ & weight $w(n,k)$ & held out\\
\midrule
$\gamma$ & 3 & $\frac13\HH3n-\frac16\HH3k$ & $n=60..72$\\
$\mathbf A$ & 2 & $\frac14\HH2k+\frac34H_k(H_k-H_{n-k})$ & $n=60..72$\\
$\mathbf D$ & 2 & $\frac15\bigl[\HH2n+H_k(2H_k-H_{n-k}-H_n)\bigr]$ & $n=60..72$\\
$s_{10}$ & 2 & $\frac15\HH2k+\frac45H_k(H_k-H_{n-k})$ & $n=60..72$\\
$s_7$ & 2 & 8 terms in $H^{(1,2)}$ at $\{k,n-k,n,2k\}$ & $n=60..72$\\
$\varepsilon$ & 3 & 10 terms in $H^{(1,2,3)}$ at $\{k,n-k,2k,2k-n\}$ & $n=60..72$\\
$\alpha$ (Domb) & 3 & 14 terms in $H^{(1,2,3)}$ at $\{k,n-k,2k,2n-2k\}$ & $n=60..72$\\
$\mathbf E$ & 2 & $\frac12K^{(2)}_{2k}+\frac34H_k(K_{2k}-K_{2n-2k})-\frac12H_{2k}(K_{2k}-K_{2n-2k})$, $K=K_{\chi_{-4}}$ & $n=70..82$\\
\bottomrule
\end{tabular}
\end{table}

The $\gamma$ line, multiplied by $6$, is the minimal closed form
\eqref{eq:minimal-intro}; Section~\ref{sec:minimal} proves it. Two structural comments.

\begin{remark}[a uniform shape for $\sum_k\bin nk^d$]\label{rem:samefamily}
$\mathbf A$ and $s_{10}$ have literally the same shape,
$\frac1{d+1}\HH2k+\frac d{d+1}H_k(H_k-H_{n-k})$ for $A^{(d)}(n)=\sum_k\bin nk^d$ with
$d=3,4$. This is the natural source of Chamberland--Straub's conjecture
\cite[eq.~(7.4)]{ChamStraub} that the Ap\'ery limit of $A^{(d)}$ is $\zeta(2)/(d+1)$, and
suggests it for every $d$ \VerifiedR{$d=3,4$ only; the recurrences for $d\ge5$ have order
$\ge3$ and were not tested}.
\end{remark}

\begin{observation}[negative results, and the dichotomy]\label{obs:negatives}
\Verified\ Exact inconsistency of the fitting systems:
\begin{itemize}[leftmargin=1.4em]
\item $\mathbf C$: no decomposition on the summand $\bin nk^2\bin{2k}k$ exists in
  \emph{any} basis tried --- plain harmonic monomials at $\{n,k,n-k,2k\}$ (14 elements),
  at $\{n,k,n-k,2k,n+k,2n-2k,2n,2n-k\}$ (44 elements), and with $\chi_{-3}$-letters
  adjoined (44 and 152 elements, 110 and 220 exact equations). All inconsistent.
\item $\mathbf E$: inconsistent in every \emph{pure} harmonic basis tried (up to 44
  elements, 8 arguments), and consistent the moment $\chi_{-4}$-letters are allowed. This
  is the decisive experiment.
\item $\delta$: inconsistent in the weight-3 basis at $\{n,k,3k,n-3k,n+k\}$ (65 elements,
  110 equations).
\item $\mathbf B$, $\mathbf F$, $\zeta$, $\eta$, $s_{18}$: not attempted (double or
  alternating sums; $\zeta$ needs a genuine two-index ansatz).
\end{itemize}
\end{observation}

The dichotomy behind this is elementary but decisive. A pure harmonic monomial
$\prod_t\HH{r_t}{x_t}$ can only ever produce values in the ring generated by
$\zeta(2),\zeta(3),\dots$ in the limit; it can never produce $L(\chi,w)$. Hence \emph{no}
$\chi$-twisted sequence can have a pure-$\zeta$ harmonic decomposition, and the letters
must be twisted --- confirmed by $\mathbf E$, and the reason $\mathbf C$'s plain fits are
inconsistent. $\mathbf C$'s residual obstruction is that its conductor-$3$ letters need
arguments that $\bin nk^2\bin{2k}k$ does not provide; the right move there is a
constant-term representation, not a bigger monomial basis \cite{Gorodetsky}. This remains
\Open.

\subsection{Verdict: which of the fifteen are proved}
\label{ssec:verdict}

We now state precisely which families Theorem~\ref{thm:LB} covers. \emph{The labels below
are the paper's evidence classes and are not to be read as claims of uniform coverage.}

\paragraph{\Proved\ --- hypotheses \textup{(H1)--(H5)} verified, Theorem~\ref{thm:LB}
applies.}

\begin{center}\small
\begin{tabular}{@{}llll@{}}
\toprule
seq. & $\chi$ & primes & why (H4) holds\\
\midrule
$\gamma$ (Ap\'ery $\zeta(3)$) & $1$ & $p\ge5$ & arguments $\{n,k\}$; coefficients $\frac13,\frac16$\\
$\mathbf A$ (Franel) & $1$ & $p\ge5$ & arguments $\{k,n-k\}$; coefficients $\frac14,\frac34$\\
$\mathbf D$ (Ap\'ery $\zeta(2)$)$^{\dagger}$ & $1$ & $p\ge7$ & arguments $\{n,k,n-k\}$; coefficients $\frac15,\frac25$\\
$s_{10}$ ($\sum_k\bin nk^4$) & $1$ & $p\ge7$ & arguments $\{k,n-k\}$; coefficients $\frac15,\frac45$\\
\bottomrule
\end{tabular}
\end{center}

%% FIXED [MAJOR-1]: D is the unique defective row. Theorem \ref{thm:LB} controls the explicit
%% FIXED sum B; identifying that sum with D's actual second solution needs (LB)=0, which this
%% FIXED paper labels \VerifiedR{n<=25} at sec-minimal.tex and whose proof by the paper's own
%% FIXED template is obstructed by Proposition \ref{prop:obstruction}. gamma, A and s_10 do not
%% FIXED have this problem: for them (LB)=0 is \Proved\ in Section \ref{sec:minimal}.
\noindent
$^{\dagger}$\emph{One qualification, and it applies to $\mathbf D$ alone.} What
Theorem~\ref{thm:LB} controls is the explicit harmonic sum $B$. For $\gamma$, $\mathbf A$ and
$s_{10}$ the identity $(LB)=0$ --- which says that this $B$ \emph{is} the sequence's second
solution --- is \Proved\ in Section~\ref{sec:minimal}. For $\mathbf D$ it is not: it is
\VerifiedR{$n\le25$} only, and Proposition~\ref{prop:obstruction} shows that the template used
for the other three provably cannot supply it. So the $\mathbf D$ row is \Proved\ as a statement
about the explicit sum $B$, and \Verified\ as a statement about the $\zeta(2)$-Ap\'ery second
solution.

\noindent
For these, (H1)--(H3) hold by the argument of Section~\ref{sec:weight3}: for
$S=\bin nk^{\alpha}\bin{n+k}n^{\beta}$ the surviving set is $\{s\le r\}$ (no borrow in
$\bin nk$) intersected with $\{r+s<p\}$ (no carry in $\bin{n+k}n$), a product region;
there $\bin nk\equiv\bin ab\bin rs$ and $\bin{n+k}n\equiv\bin{a+b}a\bin{r+s}r$ by
Lemma~\ref{lem:lucas-step}, and the deleted $r+s\ge p$ terms of $A(r)$ vanish modulo $p$
--- this is verbatim Lemma~\ref{lem:fibre}. Digit compatibility holds because
$\lfloor k/p\rfloor=b$, $\lfloor(n-k)/p\rfloor=a-b$ (no borrow) and
$\lfloor n/p\rfloor=a$.

The two steps of the proof of Theorem~\ref{thm:LB} --- ``the vanishing layer is
$\equiv0$'' and ``the surviving layer is $\equiv\chi(p)^eB_aA_r$'' --- were also checked
separately and exactly on all cells $n=ap+r<p^2$: zero failures for $\gamma$, $\mathbf A$,
$s_{10}$ at $p=5,7,11,13$ (both layers); zero failures for $\mathbf D$ at $p=7,11,13$, but
$11$ failing cells at $p=5$, which is exactly the (H4) coefficient condition biting
($\mathbf D$'s coefficients are $\frac15(1,2,-1,-1)$). So $\mathbf D$'s proof is
unconditional for $p\ge7$, and $p=5$ is \Verified\ only. $s_{10}$'s coefficients are also
$\frac15(1,4,-4)$ yet it exhibits no defect at $p=5$ --- the layer valuations happen to be
one higher there --- but the proof as written still needs $p\ge7$ for $s_{10}$. No other
defect of any kind appears in the tame cases: the only obstruction the mechanism ever
meets is the (H4) coefficient condition.

\paragraph{Need the Lemma-D upgrade --- decomposition known, tameness fails.}
\label{ssec:lemmaD}
For $\alpha$ (Domb), $\varepsilon$, $s_7$ and $\mathbf E$ a decomposition is known
(Table~\ref{tab:decomp}), but (H4) fails: the arguments reach $2n$ (letters at $2k$,
$2n-2k$, $2k-n$), so $p^ww(n,k)$ acquires a pole on the vanishing layer and the two-layer
split of the proof of Theorem~\ref{thm:LB} is \emph{invalid termwise}. \Verified: the
vanishing-layer and surviving-layer defects occur in exactly equal numbers and cancel,
which is why the congruence is nonetheless true. What is needed is a valuation bound
showing that a Kummer carry in the wide binomial ($\bin{2k}k$, $\bin{2k}n$) beats the
order-one pole of $H^{(r)}(\lfloor2k/p\rfloor)$. ($s_7$ additionally excludes $p=7$: its
coefficients have denominator $28$; \Verified\ $13$ failing cells at $p=7$ and none at
$p=5,11,13$.) Notably $\mathbf E$ is the only known instance of the twisted theorem
($e=1$), and it is non-tame: attempts to find a \emph{tame} $\chi_{-4}$-decomposition for
$\mathbf E$ (arguments $\{n,k,n-k\}$, and $\{n,k,n-k,2k\}$) were \VerifiedR{inconsistent}
--- the wide arguments are essential. So the twisted case of Theorem~\ref{thm:LB} is
proved as a mechanism but has, at present, no tame instance. This is \Open.

\paragraph{Remain conjectural.}
For $\mathbf B$, $\mathbf C$, $\mathbf F$, $\delta$, $\zeta$, $\eta$ and $s_{18}$ no
harmonic decomposition is known, and Conjecture~\ref{conj:master} is \Verified\ only. Of
these, $\mathbf C$ and $\delta$ carry \emph{proved-negative} results in the natural bases
(Observation~\ref{obs:negatives}), so they need genuinely new letters: for $\mathbf C$ the
conductor-$3$ letters, for $\delta$ presumably letters attached to the $3$-section of the
summand.

\subsection{The sequences with no archimedean Ap\'ery limit}
\label{ssec:nolimit}

\begin{observation}\label{obs:nolimit}
\Verified\ Three of the fifteen have no Ap\'ery limit at all. For \eqref{eq:R2} the
characteristic roots solve $\lambda^2-a\lambda+c=0$, and for \eqref{eq:R3},
$\lambda^2-2a\lambda+c=0$. The discriminants are
\[
   \mathbf B:\;-27,\qquad \delta:\;-128,\qquad \eta:\;-16,
\]
i.e.\ complex-conjugate pairs of equal modulus, so Poincar\'e--Perron provides no dominant
solution and $B(n)/A(n)$ does not converge --- numerically, $0$ digits of agreement
between $n=590$ and $n=600$, against $380$ digits for the convergent sequences. Every
other sporadic sequence has real distinct roots.
\end{observation}

Yet all three satisfy Conjecture~\ref{conj:master} with a flat floor, with
$\chi=\chi_{-3}$, trivial, and $\chi_5$ respectively. \emph{The $p$-adic Ap\'ery limit
exists where the archimedean one does not.} This is the cleanest evidence available to us
that the descent is a statement about Frobenius and not about the archimedean limit; we
return to it in Section~\ref{sec:discussion}.

\section{Minimal closed forms for second solutions}
\label{sec:minimal}

Theorem~\ref{thm:LB} needs a harmonic-monomial weight. For Ap\'ery's $b_n$ the fitting
procedure of Section~\ref{ssec:decomp} returned one of startling brevity, and this section
proves that it is correct. We then prove two more closed forms of the same kind, and
record a sharp obstruction: the method provably fails for the $\zeta(2)$-Ap\'ery
sequence.

\subsection{The minimal form}

\begin{theorem}[minimal form]\label{thm:minimal}
\Proved\ For every integer $n\ge0$,
\[
   \sum_{k=0}^{n}\bin nk^2\bin{n+k}k^2\bigl(2\HH3n-\HH3k\bigr)
   \;=\;
   \sum_{k=0}^{n}\bin nk^2\bin{n+k}k^2
     \Bigl(\HH3n+\sum_{m=1}^{k}\frac{(-1)^{m-1}}{2m^3\bin nm\bin{n+m}m}\Bigr)
   \;=\;b_n .
\]
Equivalently,
$\sum_{k=0}^{n}A(n,k)\bigl(\HH3n-\HH3k-\sum_{m=1}^ku(n,m)\bigr)=0$.
\end{theorem}

Write $B_{\min}(n)$ for the left-hand sum. The proof is by the classical route: both
sides satisfy Ap\'ery's recurrence and agree at $n=0,1$. Set
\begin{equation}\label{eq:Ldef}
   P(n):=34n^3+51n^2+27n+5,\qquad
   (Lu)(n):=(n+1)^3u(n+1)-P(n)u(n)+n^3u(n-1),
\end{equation}
and, for integers $n\ge1$, $k\ge0$, the \emph{base term}
\begin{equation}\label{eq:Phi}
   \Phi(n,k):=\frac{\bin{n+1}k^2\bin{n+k-1}k^2}{n^2(n+1)^2}.
\end{equation}

\begin{lemma}[conversion]\label{lem:conv}
\Proved\ For all integers $n\ge1$, $k\ge0$:
\begin{align}
   A(n,k)&=\Phi(n,k)(n+1-k)^2(n+k)^2, \tag{0a}\\
   A(n+1,k)&=\Phi(n,k)(n+k)^2(n+k+1)^2, \tag{0b}\\
   A(n-1,k)&=\Phi(n,k)(n+1-k)^2(n-k)^2, \tag{0c}\\
   \Phi(n,k+1)(k+1)^4&=\Phi(n,k)(n+1-k)^2(n+k)^2. \tag{0d}
\end{align}
\end{lemma}

\begin{proof}
Each reduces to the binomial conversions
$\bin nk(n+1)=\bin{n+1}k(n+1-k)$, $\bin{n+k}kn=\bin{n+k-1}k(n+k)$,
$\bin{n+k+1}kn(n+1)=\bin{n+k-1}k(n+k)(n+k+1)$,
$\bin{n-1}kn(n+1)=\bin{n+1}k(n+1-k)(n-k)$, and, for (0d),
%% FIXED [MINOR]: only the last two are literal instances of $\bin N{j+1}(j+1)=\bin Nj(N-j)$;
%% FIXED the first four move the UPPER index and are instances of the companion form. The
%% FIXED hedge present at MINIMAL_FORM_PROOF.md:48-51 had been dropped.
$\bin{n+1}{k+1}(k+1)=\bin{n+1}k(n+1-k)$, $\bin{n+k}{k+1}(k+1)=\bin{n+k-1}k(n+k)$. The last
two, those of (0d), are absorption identities $\bin N{j+1}(j+1)=\bin Nj(N-j)$ --- the second of
them followed by one step of the companion. The first four move the \emph{upper} index instead
and are instances of that companion, $\bin{N+1}j(N+1-j)=\bin Nj(N+1)$, applied once or twice.
All six hold for
\emph{every} integer $k\ge0$, including the degenerate range $k>n$ where both sides
vanish: the conversions are stated multiplicatively, so no $0/0$ arises.
\end{proof}

Lemma~\ref{lem:conv} is the device that removes every boundary subtlety: after it, all
identities below are $\Phi(n,k)$ times polynomial identities in $\Q[n,k]$.
\VerifiedR{all six conversions, $n=1..12$, $k=0..n+4$, difference exactly $0$}.

\subsection{Step B: the Zeilberger certificate for \texorpdfstring{$a_n$}{a\_n}}

Set
\begin{equation}\label{eq:TG}
   T(n,k):=4(2n+1)\bigl(2k^2-3k-4n^2-4n\bigr),
   \qquad
   G(n,k):=\Phi(n,k)\,k^4\,T(n,k).
\end{equation}

\begin{lemma}[CERT-1]\label{lem:cert1}
\Proved\ As an identity in $\Q[n,k]$,
\begin{multline*}
  T(n,k+1)(n+1-k)^2(n+k)^2-k^4T(n,k)\\
  =(n+1)^3(n+k+1)^2(n+k)^2-P(n)(n+1-k)^2(n+k)^2+n^3(n-k)^2(n+1-k)^2 .
\end{multline*}
\end{lemma}

\begin{proposition}\label{prop:B}
\Proved\ For all integers $n\ge1$, $k\ge0$,
\begin{equation}\label{eq:dagger}
   (n+1)^3A(n+1,k)-P(n)A(n,k)+n^3A(n-1,k)\;=\;G(n,k+1)-G(n,k).
\end{equation}
\end{proposition}

\begin{proof}
By (0a)--(0c) the left side is $\Phi(n,k)$ times the right-hand side of
Lemma~\ref{lem:cert1}. By (0d),
$G(n,k+1)=\Phi(n,k+1)(k+1)^4T(n,k+1)=\Phi(n,k)(n+1-k)^2(n+k)^2T(n,k+1)$ and
$G(n,k)=\Phi(n,k)k^4T(n,k)$, so the right side is $\Phi(n,k)$ times the left-hand side of
Lemma~\ref{lem:cert1}.
\end{proof}

The boundary values are $G(n,0)=0$ (factor $k^4$), $G(n,k)=0$ for $k\ge n+2$ (factor
$\bin{n+1}k^2$), and $G(n,n+1)=-4(n+1)^3\bin{2n+1}{n+1}^2$. Summing \eqref{eq:dagger} over
$k=0,\dots,N$ with $N\ge n+1$ telescopes to $G(n,N+1)-G(n,0)=0$, which recovers Ap\'ery's
recurrence $(La)(n)=0$.

\subsection{Step B\texorpdfstring{$'$}{'} and Step C: \texorpdfstring{$L$ annihilates $B_{\min}$}{L annihilates B\_min}}

Write $w(n,k):=2\HH3n-\HH3k$, so $B_{\min}(n)=\sum_{k\ge0}A(n,k)w(n,k)$ (terms with $k>n$
vanish). Since $w(n+1,k)=w(n,k)+2/(n+1)^3$ and $w(n-1,k)=w(n,k)-2/n^3$,
\[
   (LB_{\min})(n)=\sum_{k\ge0}w(n,k)\bigl(G(n,k+1)-G(n,k)\bigr)+2a_{n+1}-2a_{n-1},
\]
by Proposition~\ref{prop:B}. Abel summation, with $w(n,k+1)-w(n,k)=-1/(k+1)^3$ and the
boundary values of $G$, turns this into
\begin{equation}\label{eq:Bprime}
   (LB_{\min})(n)\;=\;\sum_{j\ge1}\frac{G(n,j)}{j^3}+2a_{n+1}-2a_{n-1}.
\end{equation}

Now set
\begin{equation}\label{eq:UK}
   U(n,k):=-4(k-1)(2n+1)\bigl(2n^2+2n+1-k\bigr),
   \qquad
   K(n,k):=\frac{\Phi(n,k)k^4U(n,k)}{n^2(n+1)^2}.
\end{equation}

\begin{lemma}[CERT-2]\label{lem:cert2}
\Proved\ As an identity in $\Q[n,k]$,
\begin{multline*}
   U(n,k+1)(n+1-k)^2(n+k)^2-k^4U(n,k)\\
   =n^2(n+1)^2\Bigl[kT(n,k)+2\bigl((n+k)^2(n+k+1)^2-(n+1-k)^2(n-k)^2\bigr)\Bigr],
\end{multline*}
and the bracket equals $4k(2n+1)(4k^2-3k-2n-2n^2)$.
\end{lemma}

\begin{proposition}\label{prop:C}
\Proved\ For all integers $n\ge1$, $k\ge0$, with the convention $G(n,0)/0^3:=0$
(consistent, since $G(n,k)/k^3=\Phi(n,k)kT(n,k)$),
\begin{equation}\label{eq:ddagger}
   \frac{G(n,k)}{k^3}+2A(n+1,k)-2A(n-1,k)\;=\;K(n,k+1)-K(n,k).
\end{equation}
Consequently $\sum_{j\ge1}G(n,j)/j^3+2a_{n+1}-2a_{n-1}=0$, since $K(n,0)=0$ and
$K(n,k)=0$ for $k\ge n+2$.
\end{proposition}

\begin{proof}
By (0b), (0c) the left side is $\Phi(n,k)\bigl[kT(n,k)
+2((n+k)^2(n+k+1)^2-(n+1-k)^2(n-k)^2)\bigr]$; by (0d) the right side is
$\Phi(n,k)\bigl[(n+1-k)^2(n+k)^2U(n,k+1)-k^4U(n,k)\bigr]/(n^2(n+1)^2)$. Equality is
Lemma~\ref{lem:cert2}. Summing over $k=0,\dots,N$ with $N\ge n+1$ telescopes.
\end{proof}

\begin{theorem}\label{thm:LBmin}
\Proved\ $(LB_{\min})(n)=0$ for all $n\ge1$.
\end{theorem}

\begin{proof}
Combine \eqref{eq:Bprime} with Proposition~\ref{prop:C}.
\end{proof}

\subsection{Step D: the classical companion satisfies the same recurrence}

This is the half in which Ap\'ery's harmonic weight $c(n,k)$ must be differenced
\emph{in $n$}. Everything collapses because of two one-line Gosper certificates whose
entire content is the pair of polynomial identities
\begin{equation}\label{eq:miracle}
   m^2+(n+m)(n-m)=n^2,
   \qquad
   m^2+(n+1-m)(n+1+m)=(n+1)^2 .
\end{equation}

Put $D(n,k):=\bin nk\bin{n+k}k$ (so $A=D^2$), and
\[
 \begin{gathered}
   \varphi(n,m):=m\,u(n,m)=\frac{(-1)^{m-1}((m-1)!)^2(n-m)!}{2(n+m)!},\\
   \Lambda(n,k):=(n-k)\varphi(n,k+1)=\frac{(-1)^k}{2(n+k+1)D(n,k)} .
 \end{gathered}
\]
Since $D(n,k)\ne0$ for $0\le k\le n$, every division below is legitimate on its stated
range.

\begin{lemma}[the two inner Gosper certificates]\label{lem:D1}
\Proved\ For $1\le m\le n$ put
\[
   Y_1(m):=-\frac{(n+m)\varphi(n,m)}{n^2},
   \qquad
   Y_2(m):=-\frac{(n+1-m)\varphi(n,m)}{(n+1)^2} .
\]
Then, using $\varphi(n,m+1)/\varphi(n,m)=-m^2/\bigl((n-m)(n+m+1)\bigr)$,
\[
   Y_1(m+1)-Y_1(m)=\frac{\varphi(n,m)}{n-m}\ (1\le m\le n-1),
   \quad
   Y_2(m+1)-Y_2(m)=\frac{\varphi(n,m)}{n+m+1}\ (1\le m\le n),
\]
the two brackets collapsing by the identities \eqref{eq:miracle}.
\end{lemma}

\begin{lemma}[the $n$-differences of Ap\'ery's weight]\label{lem:D2}
\Proved\ Using $u(n+1,m)=u(n,m)\frac{n+1-m}{n+m+1}$ and $u(n-1,m)=u(n,m)\frac{n+m}{n-m}$,
i.e.\ $u(n+1,m)-u(n,m)=-2\varphi(n,m)/(n+m+1)$ and
$u(n-1,m)-u(n,m)=2\varphi(n,m)/(n-m)$, Lemma~\ref{lem:D1} telescopes with
$Y_2(1)=-\tfrac1{2(n+1)^3}$, $Y_1(1)=-\tfrac1{2n^3}$ to give
\[
 \begin{gathered}
   c(n+1,k)-c(n,k)=\frac{2\Lambda(n,k)}{(n+1)^2}\quad(0\le k\le n),\\
   c(n-1,k)-c(n,k)=-\frac{2(n+k+1)\Lambda(n,k)}{(n-k)n^2}\quad(0\le k\le n-1).
 \end{gathered}
\]
\end{lemma}

\begin{remark}[Ap\'ery's miracle]\label{rem:miracle}
In Lemma~\ref{lem:D2} the rational pieces $1/(n+1)^3$ and $-1/n^3$ coming from
$\HH3{n\pm1}-\HH3n$ cancel \emph{exactly} against the boundary values $Y_2(1)$, $Y_1(1)$
of the inner Gosper antidifferences, and that cancellation is precisely the polynomial
identity $m^2+(n+1-m)(n+m+1)=(n+1)^2$. This is why Ap\'ery's weight carries the $\tfrac12$,
the $(-1)^{m-1}$ and the $m^3\bin nm\bin{n+m}m$: the weight is not found, it is
\emph{forced} --- $c(n,k)$ is the unique weight, up to the kernel of the recurrence, whose
$n$-difference is a \emph{single hypergeometric term} rather than a sum.
\end{remark}

\begin{lemma}[the correction term]\label{lem:D3}
\Proved\ For $0\le k\le n$ set
$\mathcal E(n,k):=(n+1)^3A(n+1,k)\bigl(c(n+1,k)-c(n,k)\bigr)+n^3A(n-1,k)\bigl(c(n-1,k)-c(n,k)\bigr)$,
the second summand being $0$ at $k=n$ because $A(n-1,n)=0$. Then by Lemma~\ref{lem:D2},
(0b), (0c) and $A(n,k)\Lambda(n,k)=(-1)^kD(n,k)/\bigl(2(n+k+1)\bigr)$,
\[
   \mathcal E(n,k)=(-1)^kD(n,k)\Bigl[\frac{(n+1)(n+k+1)}{(n+1-k)^2}-\frac{n(n-k)}{(n+k)^2}\Bigr].
\]
\end{lemma}

\begin{proposition}\label{prop:D4}
\Proved\ For $0\le k\le n$, with $B(n-1,k):=0$ for $k\ge n$ and $B(n,k):=A(n,k)c(n,k)$,
\[
   (n+1)^3B(n+1,k)-P(n)B(n,k)+n^3B(n-1,k)
   =c(n,k)\bigl(G(n,k+1)-G(n,k)\bigr)+\mathcal E(n,k).
\]
At $k=n+1$ only the first term survives: the left side equals
$(n+1)^3A(n+1,n+1)c(n+1,n+1)$.
\end{proposition}

\begin{proof}
Substitute $c(n\pm1,k)=c(n,k)+\bigl(c(n\pm1,k)-c(n,k)\bigr)$ and use
Proposition~\ref{prop:B}.
\end{proof}

\begin{lemma}[the top term]\label{lem:D5}
\Proved\ $G(n,n+1)=-4(n+1)^3\bin{2n+1}{n+1}^2=-(n+1)^3A(n+1,n+1)$, using
$\bin{2n+2}{n+1}=2\bin{2n+1}{n+1}$. Moreover
$u(n+1,n+1)=\frac{(-1)^n}{4(n+1)^3\bin{2n+1}{n+1}}$ and
$\Lambda(n,n)=\frac{(-1)^n}{2(n+1)\bin{2n+1}{n+1}}$, so that
$c(n+1,n+1)-c(n,n)=\tfrac54\,(-1)^n\bigl/\bigl((n+1)^3\bin{2n+1}{n+1}\bigr)$ and
\[
   \mathrm{Top}(n):=c(n,n)G(n,n+1)+(n+1)^3A(n+1,n+1)c(n+1,n+1)
   =5(-1)^n\bin{2n+1}{n+1}.
\]
\end{lemma}

This is where the classical $5$ of $\zeta(3)=\tfrac52\sum_m(-1)^{m-1}m^{-3}\bin{2m}m^{-1}$
enters.

Summing Proposition~\ref{prop:D4} over $k=0,\dots,n$, adding the $k=n+1$ term, and
Abel-summing $\sum_{k=0}^nc(n,k)(G(n,k+1)-G(n,k))=c(n,n)G(n,n+1)-\sum_{j=1}^nu(n,j)G(n,j)$
(using $G(n,0)=0$), we get
\begin{equation}\label{eq:Lb-assembly}
   (Lb)(n)=\mathrm{Top}(n)+\sum_{k=0}^{n}V(n,k),
   \qquad V(n,k):=\mathcal E(n,k)-u(n,k)G(n,k),
\end{equation}
the $k=0$ term of the second piece being $0$. Since
\[
  u(n,k)G(n,k)=\frac{(-1)^{k-1}D(n,k)\,k\,T(n,k)}{2(n+1-k)^2(n+k)^2},
\]
we have $V(n,k)=(-1)^kD(n,k)W(n,k)$ with
\[
   W(n,k)=\frac{(n+1)(n+k+1)}{(n+1-k)^2}-\frac{n(n-k)}{(n+k)^2}
          +\frac{kT(n,k)}{2(n+1-k)^2(n+k)^2}.
\]

\begin{lemma}[CERT-3]\label{lem:cert3}
\Proved\ As an identity in $\Q[n,k]$,
\[
   2(n+1)(n+k+1)(n+k)^2-2n(n-k)(n+1-k)^2+kT(n,k)+10k(2n+1)(n+1-k)(n+k)=0,
\]
that is,
\[
 \begin{gathered}
   W(n,k)=\frac{-5k(2n+1)}{(n+1-k)(n+k)},\qquad\text{whence}\\
   V(n,k)=\frac{-5(2n+1)}{n(n+1)}\,(-1)^k\,k\,\bin{n+1}k\bin{n+k-1}k,
 \end{gathered}
\]
using $\bin nk/(n+1-k)=\bin{n+1}k/(n+1)$ and $\bin{n+k}k/(n+k)=\bin{n+k-1}k/n$.
\end{lemma}

The entire weight-three structure thus collapses to a single linear-over-quadratic
rational function.

\begin{lemma}[residual binomial identity]\label{lem:D6}
\Proved\
\[
   \text{(D-BIN)}\qquad
   \sum_{k=0}^{n}(-1)^k\,k\,\bin{n+1}k\bin{n+k-1}k\;=\;(-1)^n\,n\bin{2n}n .
\]
\end{lemma}

\begin{proof}
Put $s(n,k):=(-1)^kk\bin{n+1}k\bin{n+k-1}k$ and $Y(n,k):=-k(k-1)s(n,k)/(n(n+1))$. The
absorption identities $\bin{n+1}{k+1}(k+1)=\bin{n+1}k(n+1-k)$ and
$\bin{n+k}{k+1}(k+1)=\bin{n+k-1}k(n+k)$ of Lemma~\ref{lem:conv} give
$(k+1)s(n,k+1)=-(-1)^k\bin{n+1}k\bin{n+k-1}k(n+1-k)(n+k)$, hence for \emph{every} integer
$k\ge0$, division-free,
\begin{multline*}
   n(n+1)\bigl(Y(n,k+1)-Y(n,k)\bigr)\\
   =(-1)^kk\bin{n+1}k\bin{n+k-1}k\bigl[(n+1-k)(n+k)+k(k-1)\bigr]=n(n+1)s(n,k),
\end{multline*}
by the polynomial identity $(n+1-k)(n+k)+k(k-1)=n(n+1)$. So $Y(n,\cdot)$ telescopes
$s(n,\cdot)$ and $\sum_{k=0}^ns(n,k)=Y(n,n+1)-Y(n,0)=Y(n,n+1)=-s(n,n+1)
=(-1)^nn\bin{2n}n$, using $\bin{2n}{n+1}=n\bin{2n}n/(n+1)$.
\end{proof}

\begin{theorem}\label{thm:Lb}
\Proved\ $(Lb)(n)=0$ for all $n\ge1$.
\end{theorem}

\begin{proof}
By \eqref{eq:Lb-assembly}, Lemma~\ref{lem:cert3} and Lemma~\ref{lem:D6},
\begin{align*}
   \sum_{k=0}^{n}V(n,k)
   &=\frac{-5(2n+1)}{n(n+1)}(-1)^nn\bin{2n}n
    =-5(-1)^n\bin{2n}n\frac{2n+1}{n+1}\\
   &=-5(-1)^n\bin{2n+1}{n+1}=-\mathrm{Top}(n). \qedhere
\end{align*}
\end{proof}

\begin{proof}[Proof of Theorem~\ref{thm:minimal}]
Theorems~\ref{thm:LBmin} and \ref{thm:Lb} say that both sides are annihilated by $L$,
whose leading coefficient $(n+1)^3$ never vanishes for $n\ge1$. It remains to check
agreement at $n=0,1$. We have $B_{\min}(0)=A(0,0)(2\HH30-\HH30)=0=b_0$, and
$B_{\min}(1)=A(1,0)\cdot2+A(1,1)\cdot1=1\cdot2+4\cdot1=6$; on the other side
$u(1,1)=1/(2\cdot1^3\bin11\bin21)=1/4$, so $c(1,0)=1$, $c(1,1)=5/4$ and
$b_1=1\cdot1+4\cdot\tfrac54=6$. Induct.
\end{proof}

\begin{table}[ht]
\centering\small
\caption{Certificate inventory for Theorem~\ref{thm:minimal}. All are verified by exact
polynomial arithmetic (\texttt{Expand}$\to0$).}
\label{tab:certs}
\begin{tabular}{@{}llp{6.4cm}@{}}
\toprule
\# & role & certificate\\
\midrule
CERT-1 & Zeilberger certificate for $a_n$ & $T(n,k)=4(2n+1)(2k^2-3k-4n^2-4n)$, $G=\Phi k^4T$\\
CERT-2 & Gosper certificate for the $w$-correction & $U(n,k)=-4(k-1)(2n+1)(2n^2+2n+1-k)$, $K=\Phi k^4U/(n^2(n+1)^2)$\\
CERT-3 & the $W$-collapse & $W(n,k)=-5k(2n+1)/\bigl((n+1-k)(n+k)\bigr)$\\
CERT-4 & inner Gosper certificate, $\Delta^+_n$ & $Y_2(m)=-(n+1-m)\varphi(n,m)/(n+1)^2$\\
CERT-5 & inner Gosper certificate, $\Delta^-_n$ & $Y_1(m)=-(n+m)\varphi(n,m)/n^2$\\
CERT-6 & Gosper certificate for (D-BIN) & $Y(n,k)=-k(k-1)s(n,k)/\bigl(n(n+1)\bigr)$\\
\bottomrule
\end{tabular}
\end{table}

The supporting polynomial identities, all \texttt{Expand}$\to0$, are
$m^2+(n+m)(n-m)=n^2$; $m^2+(n+1-m)(n+m+1)=(n+1)^2$; $(n+1-k)(n+k)+k(k-1)=n(n+1)$; and
CERT-1, CERT-2, CERT-3 as displayed. Numerically, $B_{\min}(n)-b_n=0$ exactly for
$n=0..50$ in exact rational arithmetic, in two independent implementations
\VerifiedR{$n\le50$}; and $B_{\min}$ reproduces the literature values
$0,6,\tfrac{351}4,\tfrac{62531}{36},\tfrac{11424695}{288},\tfrac{35441662103}{36000}$.

\begin{remark}[formalisation-friendliness]\label{rem:leanfriendly}
Nothing above uses Gamma functions, limits or $0/0$: Lemma~\ref{lem:conv} converts every
binomial ratio into a \emph{multiplicative} absorption identity valid for all integers
$k\ge0$, after which each step is a polynomial identity in $\Q[n,k]$ plus a finite
telescoping sum. The only analytic input is $\HH3{n+1}-\HH3n=1/(n+1)^3$.
\end{remark}

\subsection{Two more closed forms, and the mechanism behind them}

\begin{observation}[symmetrisation]\label{obs:symm}
\Proved\ For the $d$-fold binomial sums $A^{(d)}(n)=\sum_k\bin nk^d$ the summand is
symmetric under $k\mapsto n-k$, so the fitted weights of Table~\ref{tab:decomp} may be
replaced by their $k\leftrightarrow n-k$ averages without changing the sum:
\[
   \tfrac1{d+1}\HH2k+\tfrac d{d+1}H_k(H_k-H_{n-k})
   \;\rightsquigarrow\;
   \tfrac1{2(d+1)}\bigl(\HH2k+\HH2{n-k}\bigr)+\tfrac d{2(d+1)}\bigl(H_k-H_{n-k}\bigr)^2 .
\]
\VerifiedR{exactly, $n=0..12$, $d=3,4$}. Concretely,
\[
 \begin{gathered}
   B_{\mathbf A}(n)=\sum_k\bin nk^3\Bigl[\tfrac18\bigl(\HH2k+\HH2{n-k}\bigr)+\tfrac38\bigl(H_k-H_{n-k}\bigr)^2\Bigr],\\
   B_{s_{10}}(n)=\sum_k\bin nk^4\Bigl[\tfrac1{10}\bigl(\HH2k+\HH2{n-k}\bigr)+\tfrac25\bigl(H_k-H_{n-k}\bigr)^2\Bigr].
 \end{gathered}
\]
\end{observation}

This is not cosmetic. In the symmetric form the only weight-one letter that ever appears
is the single antisymmetric combination $\delta(n,k):=H_k-H_{n-k}$, because
$\sigma:=\HH2k+\HH2{n-k}$ differences to \emph{rational} functions under both $n\mapsto
n\pm1$ and $k\mapsto k+1$. The certificate module is therefore of rank $2$
($\{1,\delta\}$) instead of rank $4$ ($\{1,H_k,H_{n-k},H_n\}$), and one faces two Gosper
problems instead of four.

The template is then: (i) a Zeilberger certificate
$G(n,k)=S(n,k)k^d\rho(n,k)/(n+1-k)^d$, whose defining equation collapses to one polynomial
equation in $\Q[n,k]$ because $S(n,k+1)/S(n,k)=((n-k)/(k+1))^d$, so that the shifted
certificate is the summand times a polynomial; (ii) a split of the recurrence applied to
$S\cdot w$ into $w\cdot\Delta_kG$ plus a correction whose $n$-differences lie in
$\Q(n,k)+\Q(n,k)\delta$; (iii) Abel summation; (iv) two Gosper problems
\[
 \begin{gathered}
   r(n,k)z_1(n,k+1)-z_1(n,k)=B(n,k),\\
   r(n,k)z_0(n,k+1)-z_0(n,k)=A(n,k)-r(n,k)z_1(n,k+1)e(n,k),
 \end{gathered}
\]
with $r=S(n,k+1)/S(n,k)$ and $e=\frac1{k+1}+\frac1{n-k}$, both solved by rational ans\"atze
$z_1=N(k)/(n+1-k)^{p}$, $z_0=M(k)/(n+1-k)^{p+1}$; and (v) three scalar boundary
identities, of which the first,
\[
   \text{(b1)}\qquad \rho(n,n+1)=-c_+(n)\,S(n+1,n+1)/S(n,n),
\]
is structural: it says that the Zeilberger certificate at the top of its support
reproduces the leading term of the recurrence, and it is what forces a harmonic
decomposition to exist at all.

%% FIXED [MAJOR-2]: \Proved\ is defined at sec-intro.tex as "a complete proof is written out
%% FIXED here", and neither this theorem nor Theorem \ref{thm:s10} writes one out: the Gosper
%% FIXED equations the certificates solve are not stated, and Theorem \ref{thm:s10}'s N_10, M_10
%% FIXED are not in the paper at all. Both are now labelled for what they are: certificates
%% FIXED exhibited and checked, with the verification exact.
\begin{theorem}[Franel]\label{thm:franel}
\Proved\ modulo the certificate check described after Theorem~\ref{thm:s10}; the displayed
identities are \Proved\ (\texttt{Factor}$\to0$, exact).
With $A(n)=\sum_k\bin nk^3$ (A000172), $(n+1)^2u_{n+1}=(7n^2+7n+2)u_n+8n^2u_{n-1}$
and $B(n)=\sum_k\bin nk^3\bigl[\tfrac18(\HH2k+\HH2{n-k})+\tfrac38(H_k-H_{n-k})^2\bigr]$,
one has $(LB)=0$ for the Franel operator, so $B$ is the second solution, with values
$B(0),\dots,B(8)$ equal to
$0$, $1$, $4$, $\tfrac{208}9$, $\tfrac{1280}9$, $\tfrac{208384}{225}$,
$\tfrac{1404928}{225}$, $\tfrac{95174656}{2205}$, $\tfrac{3351248896}{11025}$.
The certificates are
\[
 \begin{gathered}
  \rho_F(n,k)=-\tfrac1n\bigl(4-12k+12k^2-4k^3+22n-39kn+18k^2n+32n^2-27kn^2+14n^3\bigr),\\
  G_F=\tfrac{k^3\rho_F\bin{n+1}k^3}{(n+1)^3},
 \end{gathered}
\]
$z_1=N_F(k)/(n+1-k)^4$, $z_0=M_F(k)/(n+1-k)^5$ with
{\small
\begin{align*}
  N_F&=\tfrac{3k^2}{4n}\bigl(4-16k+24k^2-16k^3+4k^4+26n-68kn+63k^2n-20k^3n\\
     &\qquad\qquad+54n^2-88kn^2+39k^2n^2+46n^3-36kn^3+14n^4\bigr),\\
  M_F&=-\tfrac{k}{2n}\bigl(2-10k+20k^2-20k^3+10k^4-2k^5+15n-50kn+70k^2n-45k^3n+11k^4n\\
     &\qquad\qquad+40n^2-90kn^2+80k^2n^2-25k^3n^2+50n^3-70kn^3+30k^2n^3\\
     &\qquad\qquad+30n^4-20kn^4+7n^5\bigr),
\end{align*}}
and the boundary identities, all \texttt{Factor}$\to0$, are
$\rho_F(n,n+1)=-(n+1)^2$, $N_F(n)=\tfrac34\bigl((n+1)^5-(n+1)\bigr)$,
$M_F(n)=-\tfrac12\bigl(1+(n+1)^5\bigr)$.
\end{theorem}

\begin{theorem}[$s_{10}=\sum_k\bin nk^4$]\label{thm:s10}
\Proved\ modulo the certificate check described below.
With $(n+1)^3u_{n+1}=(2n+1)(6n^2+6n+2)u_n+n(64n^2-4)u_{n-1}$ and
$B(n)=\sum_k\bin nk^4\bigl[\tfrac1{10}(\HH2k+\HH2{n-k})+\tfrac25(H_k-H_{n-k})^2\bigr]$,
one has $(LB)=0$, with $B(0),\dots,B(6)$ equal to
$0$, $1$, $\tfrac{21}4$, $\tfrac{1001}{18}$, $\tfrac{85085}{144}$,
$\tfrac{4203199}{600}$, $\tfrac{52055003}{600}$.
The certificate is $G_{10}=k^4\rho_{10}\bin{n+1}k^4/(n+1)^4$ with
{\small
\begin{align*}
 \rho_{10}(n,k)=-\tfrac1{n^3}\bigl(&-4k+16k^2-24k^3+16k^4-4k^5\\
 &+10n-72kn+172k^2n-184k^3n+90k^4n-16k^5n\\
 &+80n^2-368kn^2+600k^2n^2-416k^3n^2+104k^4n^2\\
 &+255n^3-796kn^3+818k^2n^3-276k^3n^3\\
 &+385n^4-756kn^4+374k^2n^4+275n^5-260kn^5+75n^6\bigr),
\end{align*}}
$z_1=N_{10}(k)/(n+1-k)^5$ and $z_0=M_{10}(k)/(n+1-k)^6$ (both of degree $9$ in $k$, with
factors $k^3$ and $k^2$; full coefficients in the certificate file, see
Section~\ref{sec:data}), and boundary identities $\rho_{10}(n,n+1)=-(n+1)^3$,
$N_{10}(n)=\tfrac45\bigl((n+1)^7-(n+1)^2\bigr)$,
$M_{10}(n)=-\tfrac12\bigl((n+1)+(n+1)^7\bigr)$, all \texttt{Factor}$\to0$.
\end{theorem}

\begin{remark}[what the two theorems above do and do not contain]\label{rem:certstatus}
Both are proved by the same route as Theorem~\ref{thm:Lb}: one exhibits $\rho$, $G$,
$z_1$, $z_0$, checks $\Delta_k G=\rho S$ and the two Gosper identities as polynomial identities
in $\Q(n)[k]$, and checks the boundary identities. Every displayed identity above was verified
that way, by exact expansion to $0$. What is \emph{not} written out here is the derivation: we
do not state the Gosper key equations that $z_1,z_0$ solve, and for Theorem~\ref{thm:s10} the
polynomials $N_{10}$, $M_{10}$ themselves are not printed --- their full coefficients are in the
certificate file (Section~\ref{sec:data}). A reader wanting a self-contained proof must read
that file. We flag this because \Proved\ is defined in Section~\ref{ssec:evidence} as ``a
complete proof is written out here'', and for these two theorems the certificate half is
delegated rather than printed.
\end{remark}

\begin{observation}[the pattern]\label{obs:pattern}
With $\beta_d=\frac d{2(d+1)}$ the $\delta^2$-coefficient and $c_+(n)=(n+1)^{d-1}$, both
theorems read
\[
 \begin{gathered}
   \rho(n,n+1)=-(n+1)^{d-1},\qquad
   N(n)=2\beta_d\bigl((n+1)^{2d-1}-(n+1)^{d-2}\bigr),\\
   M(n)=-\tfrac12\bigl((n+1)^{2d-1}+(n+1)^{d-3}\bigr).
 \end{gathered}
\]
It is natural to conjecture that this is the general-$d$ shape, i.e.\ that
Chamberland--Straub's conjecture that $A^{(d)}$ has Ap\'ery limit $\zeta(2)/(d+1)$ is
provable this way for every $d$ for which the sequence satisfies an order-two recurrence
\VerifiedR{$d=3,4$ only}.
\end{observation}

\subsection{A sharp obstruction: the \texorpdfstring{$\zeta(2)$}{zeta(2)}-Ap\'ery case}
\label{ssec:obstruction}

The template of the previous subsection does not merely fail to apply to the
$\zeta(2)$-Ap\'ery sequence $\mathbf D$; it provably cannot apply.

Let $A(n)=\sum_k\bin nk^2\bin{n+k}k$ (A005258: $1,3,19,147,1251,\dots$),
$(n+1)^2u_{n+1}=(11n^2+11n+3)u_n+n^2u_{n-1}$, and
$B(n)=\sum_k\bin nk^2\bin{n+k}k\cdot\frac15\bigl[\HH2n+H_k(2H_k-H_{n-k}-H_n)\bigr]$, with
$(LB)=0$ \VerifiedR{$n\le25$}. Everything up to the Gosper step goes through, and
beautifully: the Zeilberger certificate is the smallest in this paper,
\[
   \rho_{\mathbf D}(n,k)=k^2+k(1+6n)-4-15n-11n^2,
   \qquad
   G_{\mathbf D}(n,k)=\frac{k^3\rho_{\mathbf D}(n,k)\bin{n+1}k^2\bin{n+k-1}k}{n(n+1)^2},
\]
\Proved\ (\texttt{Expand}$\to0$; \VerifiedR{pointwise, $n=1..9$, $k=0..n+3$}), and the
structural boundary identity (b1) holds exactly:
\[
   \rho_{\mathbf D}(n,n+1)=-2(n+1)(2n+1)=-(n+1)^2\frac{\bin{2n+2}{n+1}}{\bin{2n}n}
   =-c_+(n)\frac{S(n+1,n+1)}{S(n,n)} .
\]
The correction decomposes as $Y(n,k)/S(n,k)=A+B_1H_k+B_2(H_{n-k}+H_n)$ with
$B_2=\rho_{\mathbf D}(n,k+1)/\bigl(5(k+1)\bigr)$ \VerifiedR{exactly, $n=2..8$,
$0\le k\le n-1$}.

\begin{proposition}[the $H_{n-k}$ component is not Gosper-summable]\label{prop:obstruction}
\Proved\ Let $T(k):=S_{\mathbf D}(n,k)B_2(n,k)$. In Gosper's decomposition
$T(k+1)/T(k)=\frac{a(k)}{b(k)}\frac{c(k+1)}{c(k)}$ one has $a=(n-k)^2(n+k+1)$,
$b=(k+1)^2(k+2)$, $c(k)=\rho_{\mathbf D}(n,k+1)$, and $\gcd\bigl(a(k),b(k+j)\bigr)=1$ for
all $j\ge0$ (the roots are $n,-n-1$ versus $-j-1,-j-2$, with $n$ generic). Gosper's key
equation is $a(k)x(k+1)-b(k-1)x(k)=c(k)$. Since
\[
   a(k)-b(k-1)=-n\bigl(k^2+(n+2)k-n(n+1)\bigr)
\]
has degree $2$ with leading coefficient $-n\ne0$, the degree bound forces $\deg x=0$; and
$x=x_0$ constant forces $x_0=-1/n$ from the $k^2$ coefficient, which then produces the
$k$-coefficient $n+2$ against the required $3+6n$. Contradiction. Independently
\Verified: the corresponding linear system has no solution for any $\deg x\le8$.
\end{proposition}

\begin{remark}[diagnosis]\label{rem:diagnosis}
In the Franel and $s_{10}$ cases the weight is $k\leftrightarrow n-k$ symmetric, the
certificate module has rank $2$, and each component telescopes. Here the weight is
genuinely three-lettered ($H_k$, $H_{n-k}$, $H_n$), the components do not telescope
separately, and the weight-two part of any enlarged ansatz is forced to vanish
($\Delta(Sz)=0\Rightarrow z=0$). Closing this case needs either a different gauge for $w$
--- the harmonic-monomial fit has a large kernel, $11$-dimensional in the $\zeta(3)$ case,
and a kernel shift changes the antidifference --- or a genuine $\Pi\Sigma$-field ansatz
allowing nested sums rather than rational $\times$ harmonic-monomial. This is \Open, and
it is the first place in this programme where the rank-two mechanism is provably
insufficient.
\end{remark}

\subsection{Literature}
\label{ssec:minlit}

The recurrence $(Lb)=0$ itself is not new. It is proved in full by Schneider
\cite{Schneider2007} using Karr's $\Pi\Sigma$ difference fields and creative
telescoping, with a certificate whose numerator, after the substitution $n\mapsto n-1$,
is exactly the numerator of our $G(n,k)$ --- an independent confirmation of the
normalisation of \eqref{eq:TG}. Historically it goes back to the Zagier--Cohen
reorganisation of Ap\'ery's proof reported by van der Poorten
\cite[\S8]{vdPoorten1979}, who gives
$B_{n,k}=4(2n+1)\bigl(k(2k+1)-(2n+1)^2\bigr)\bin nk^2\bin{n+k}k^2$ and, for the second
solution, $A_{n,k}=B_{n,k}c_{n,k}+\bigl[5(2n+1)(-1)^{k-1}k/(n(n+1))\bigr]\bin nk\bin{n+k}k$
--- precisely our Step-D certificate, with $G(n,k)=B_{n,k-1}$ and $V$ as in
Lemma~\ref{lem:cert3}. But van der Poorten dispatches the verification as ``after some
massive reorganisation (9) becomes $A_{n,k}-A_{n,k-1}$'': it is exhibited, not written
out; Fischler \cite[\S1.2]{Fischler2004} likewise calls it ``une simple
v\'erification''. Zudilin \cite[eqs.~(5)--(6)]{Zudilin0202159} has the rational-function
form $s_n(t)=4(2n+1)(-2t^2+t+(2n+1)^2)$, the same polynomial under $t\leftrightarrow-k$.

What appears to be new is the \emph{minimal form} itself. We could not locate
$2\HH3n-\HH3k$ anywhere: not in OEIS \texttt{A059415} (which is exactly this $b_n$, and
whose only recorded closed form is the classical double sum), nor in
\cite{ChamStraub,Gorodetsky,Straub1401,Cooper2302}, nor in
\cite{vdPoorten1979,Schneider2007,Fischler2004,
PauleSchneider2003}. The nearest published relative is Nesterenko's Lemma~1
\cite{Nesterenko1996}, reproduced as formula (3) of \cite{Zudilin0202159},
which amounts to the weight $\HH3k+(2H_k-H_{n+k}-H_{n-k})\HH2k$ --- a longer
representation, not the minimal one. (The harmonic-monomial fitting system for $b_n$ has
an $11$-dimensional kernel, so such representations are far from unique; the pivot choice
delivered the shortest.) Similarly, (D-BIN) of Lemma~\ref{lem:D6} we could not locate as a
named identity, though it is one line from Chu--Vandermonde: from
$\sum_{k\ge0}(-1)^k\bin{n+1}k\bin{n+k-1}kx^k={}_2F_1(-(n+1),n;1;x)$, apply $x\,d/dx$ at
$x=1$ to get $-n(n+1)\,{}_2F_1(-n,n+1;2;1)=0$ for $n\ge1$, and peel off the $k=n+1$ term.
Its relative $\sum_k(-1)^k\bin nk\bin{n+k}k=(-1)^n$ is $P_n(-1)$. Finally, Paule and
Schneider \cite{PauleSchneider2003} prove Ahlgren's
$\sum_j(1-4jH_j+4jH_{n-j})\bin nj^4=(-1)^n\bin{2n}n$ --- the same right-hand side and the
same $d=4$ summand as Theorem~\ref{thm:s10} --- but only weight-one identities;
$H^{(3)}$ never appears there, so Theorem~\ref{thm:s10} is not covered by that work.

\section{Sharpness of the hypothesis \texorpdfstring{$p\ge5$}{p >= 5}}
\label{sec:sharpness}

Every statement in this paper carries the hypothesis $p\ge5$. This section records where
that hypothesis is used, where it can be removed, and what actually happens at $p=2,3$.

\subsection{Where \texorpdfstring{$p\ge5$}{p >= 5} enters the proofs}

In Theorem~\ref{thm:brow} the hypothesis is used twice, and only twice: the coefficient
$1/(2m^3)$ of Ap\'ery's letter $u(n,m)$ requires $p\nmid2$, and the identity
$(-1)^{jp-1}=(-1)^{j-1}$ in the proof of Lemma~\ref{lem:Wpart} requires $p$ odd. Nothing
else in Section~\ref{sec:weight3} sees the size of $p$; in particular
Theorem~\ref{thm:arow} is uniform in $p$.

In Theorem~\ref{thm:LB} the hypothesis enters through (H4), in the requirement
$c_j\in\Zp$ on the coefficients of the weight. This is not a formality: it is the only
obstruction the mechanism ever meets in the tame cases, and it bites exactly where the
denominators say it should. For $\mathbf D$, whose coefficients are $\frac15(1,2,-1,-1)$,
the two-layer split has $11$ failing cells at $p=5$ and none at $p=7,11,13$ \Verified;
for $s_7$, whose coefficients have denominator $28$, it has $13$ failing cells at $p=7$
and none at $p=5,11,13$ \Verified. The proof is therefore unconditional for $p\ge7$ in
those two cases, and $p=5$ (resp.\ $p=7$) is \Verified\ only. Curiously $s_{10}$, whose
coefficients are also $\frac15(1,4,-4)$, exhibits no defect at $p=5$ --- the layer
valuations happen to be one higher there --- but the proof as written still requires
$p\ge7$.

\subsection{Where it can be removed}

%% FIXED [MINOR x2]: (i) Theorem \ref{thm:LB} is stated for $p\ge5$, so it cannot itself be
%% FIXED "applied" at $p=2,3$; what survives there is its proof, with (H4) vacuous. (ii) the
%% FIXED Franel $p\ge5$ / $p\ne2$ mismatch now carries its cross-reference.
The minimal form of Theorem~\ref{thm:minimal} has coefficients $2$ and $-1$, both
integers, so (H4)'s coefficient condition becomes vacuous. Theorem~\ref{thm:LB} is stated for
$p\ge5$, but that hypothesis is used \emph{only} through (H4): inspecting its proof, every
other step is uniform in $p$. So the proof --- not the theorem as stated --- yields the
$\gamma$ instance for \emph{every} prime, including $2$ and $3$. This is
machine-verified: the Lean statement of Section~\ref{sec:lean} carries no hypothesis
$p\ge5$ and no hypothesis $a\ge1$. The same inspection gives the Franel instance for every
$p\ne2$ (its coefficient denominators are $4$), so that the $p\ge5$ recorded for
$\mathbf A$ in Section~\ref{sec:families} is not sharp: $p=3$ is fine.

So the exclusion of small primes is a property of the \emph{normalisation of the weight},
not of the congruence. In the natural normalisation $B(0)=0$, $B(1)=1$ the $\gamma$
coefficients are $\frac13,\frac16$ and $p\ge5$ is required; in the classical
normalisation $b_n=6B(n)$ they are $2,-1$ and it is not.

\subsection{What actually happens at \texorpdfstring{$p=2$ and $p=3$}{p = 2 and p = 3}}

The master law itself, however, does fail at the small primes, and not only for
normalisation reasons.

\begin{observation}\label{obs:small-primes}
\Verified\ With the twisted law of Conjecture~\ref{conj:master} and $n\le120$, the floor
drops below $w$ for
\[
 \begin{gathered}
   \mathbf A\;(p=2,\ \text{floor }1),\quad
   \mathbf F\;(p=2,\ 1),\quad
   \alpha\;(p=2,\ 2;\ p=3,\ 2),\\
   \gamma\;(p=3,\ 2),\quad
   \delta\;(p=3,\ 1),\quad
   \varepsilon\;(p=2,\ 2).
 \end{gathered}
\]
Independently, the iterated two-digit form $p^{2w}b_n\equiv b_{n_2}a_{n_1}a_{n_0}$ holds
at $p=5$, $w=3$, $n=25..70$ with minimum depth $1$, and fails at $p=2$ and $p=3$ with
minimum orders $-6$ and $0$ respectively.
\end{observation}

Six of the fifteen families therefore violate the master floor at $2$ or at $3$, and the
violations are of the same kind as the low-weight ratio failures of
Section~\ref{ssec:panel}: the floor does not merely dip, it drops by a full unit or more.
The boundary $p\ge5$ that this forces coincides with the boundary already familiar from
the denominator-conjecture side of this programme, where the primes $2$ and $3$ are
likewise excluded.

We record honestly that we have \emph{no} conceptual explanation of the $p=2,3$ failures.
The proof-side obstruction (coefficient denominators) accounts for individual cells, and
it is removable by renormalisation; the observed collapse of the master floor at $2$ and
$3$ for six families is a separate phenomenon, and it is \Open.

\section{Discussion}
\label{sec:discussion}

\subsection{Relation to Beukers--Vlasenko's Conjecture 7.5}
\label{ssec:BV}

The natural home for a congruence like Conjecture~\ref{conj:master} is the theory of
Dwork crystals, and it is worth being precise about how the two fit together. We read
\cite{BVthree} from its source; its \S7 treats \emph{completely symmetric Calabi--Yau
families}, and its numbered items there are 7.1 (definition), 7.2 (lemma), 7.3 (theorem),
7.4 (remark) and 7.5 (conjecture).

The setting is $f(x)=1-tg(x)$ and $F(t)=\sum_nf_nt^n$ with $f_n=\ct\bigl(g(x)^n\bigr)$;
for the sporadic families these $f_n$ are exactly our $A(n)$. Writing $F_N(t)$ for the
truncation of $F$ at $t^N$, and $G(t)$ for the series appearing in the second solution
$F(t)\log t+G(t)$ of the (second-order) Picard--Fuchs operator, \emph{the coefficients of
$G$ are our $B(n)$}, up to normalisation. What \cite{BVthree} proves --- their eq.~(7.6),
from \cite[eq.~(7)]{BVtwo} --- is that for all $m,s\ge1$,
\begin{equation}\label{eq:BV76}
   \frac{F(t)}{F(t^\sigma)}\;\equiv\;\frac{F_{mp^s}(t)}{F_{mp^{s-1}}(t^\sigma)}\pmod{p^s}.
\end{equation}
Reading off coefficients with the naive lift $t^\sigma=t^p$ and $s=1$ gives the
Dwork--Mellit--Vlasenko congruence whose sporadic specialisation is precisely the
first-row Lucas congruence $A(ap+r)\equiv A(a)A(r)\pmod p$ of
Theorem~\ref{thm:arow}. Their Conjecture~7.5 asserts that \eqref{eq:BV76} holds modulo
$p^{2s}$ --- for $m=2$ in the hypercubic and hyperoctahedral families and $m=n+1$ in the
simplicial ones --- where $\sigma$ is the \emph{excellent} Frobenius lift, the unique lift
with $q^\sigma=\gamma^{p-1}q^p$ for the canonical (mirror) coordinate
$q(t)=t\exp\bigl(G(t)/F(t)\bigr)$; by their Theorem~7.3 that lift is characterised by the
vanishing of the off-diagonal Frobenius entry. The authors state that they could not prove
7.5 for a single $g$.

Four points relate this to Conjecture~\ref{conj:master}.

\begin{enumerate}[leftmargin=1.7em]
\item \emph{They are the two rows of one $2\times2$ Frobenius matrix.} The crystal
  $\mathrm{CY}(g)(2)$ of \cite{BVthree} is free of rank two on $1/f$, $\theta(1/f)$, and
  its Frobenius has diagonal $(\lambda_0,\lambda_2)$ with $\lambda_0=F(t)/F(t^\sigma)$ the
  unit root. Conjecture 7.5 is a statement about $\lambda_0$ alone, sharpened from modulo
  $p^s$ to modulo $p^{2s}$. Conjecture~\ref{conj:master} is the statement about the
  \emph{other} diagonal entry: that Frobenius acts on the second graded piece with
  eigenvalue $\chi(p)p^{w}$. In coefficient form, their $\lambda_0$-congruence \emph{is}
  the first row; ours \emph{is} the second row. Neither formally implies the other.
\item \emph{But 7.5's standing hypothesis is a weak form of our conclusion.} The excellent
  lift exists as a lift of $\Z_p[[t]]$ precisely because
  $q(t)=t\exp(G/F)\in\Z_p[[t]]$ --- $p$-integrality of the mirror map, which is a
  statement about $G$, i.e.\ about $B(n)$. Our Corollary~\ref{cor:integrality},
  $d_n^wB_n\in\Zp$, obtained by iterating the descent over base-$p$ digits, is exactly the
  coefficient-level shadow of that integrality, and it is sharper: it does not merely
  assert integrality of the mirror map, it pins the depth at exactly $w$ and exhibits the
  eigenvalue $\chi(p)$. So the implication runs
  \[
   \begin{gathered}
    \text{Conjecture~\ref{conj:master}}\;\Longrightarrow\;
    \text{(coefficient form of mirror-map integrality)}\\
    \Longrightarrow\;
    \text{(the excellent lift is defined over $\Z_p[[t]]$)},
   \end{gathered}
  \]
  which is the standing hypothesis of Conjecture 7.5. In this precise sense our conjecture
  is \emph{upstream} of theirs, and not equivalent to it.
\item \emph{The character is invisible in the unit-root formulation and must be added.}
  Beukers and Vlasenko work with $\lambda_0=F/F^\sigma$, where no character can appear:
  the unit root of the holomorphic period is a unit, and the first-row Lucas congruence is
  untwisted for all fifteen families --- \Verified\ here and proved in
  \cite{MalikStraub}. The character surfaces only on the \emph{non-unit} eigenvalue. Any
  crystal-theoretic proof of Conjecture~\ref{conj:master} must therefore carry the
  information that, for seven of the fifteen families, the second Frobenius eigenvalue is
  $\chi(p)p^{w}$ and not $p^{w}$ --- i.e.\ that the family is attached to a quadratically
  twisted (CM) modular parametrisation. This is a concrete correction to make to any
  slogan of the form ``Dwork crystals imply Lucas congruences for the second solution''.
\item \emph{What our conjecture does not give.} Nothing towards the modulo-$p^{2s}$
  doubling that is the actual content of 7.5. Our floor is exactly $w$ and provably not
  more in the sense that equality is attained in every cell tested, so there is no hidden
  supercongruence in the second row of the kind 7.5 predicts for the first.
\end{enumerate}

We note for completeness that we could not find any second-solution version of these
congruences stated in \cite{BVthree} or in Vargas-Montoya's work
\cite{VMcanon,VMmum}; \cite{VMmum} proves Lucas congruences explicitly and exclusively for
the holomorphic solution, and \cite{VMcanon} proves $p$-integrality of $\exp(y_1/y_0)$
with no congruence attached.

\subsection{Frobenius, not the archimedean limit}
\label{ssec:frobenius}

Observation~\ref{obs:nolimit} deserves emphasis. Three sporadic sequences --- $\mathbf B$,
$\delta$, $\eta$ --- have characteristic roots forming a complex-conjugate pair of equal
modulus, so Poincar\'e--Perron gives no dominant solution and $B(n)/A(n)$ oscillates
forever: there is no Ap\'ery limit to speak of. Nonetheless all three satisfy
Conjecture~\ref{conj:master}, with a flat floor, and with characters $\chi_{-3}$, trivial
and $\chi_5$ respectively.

The $p$-adic Ap\'ery limit therefore exists where the archimedean one does not. This is
the sharpest indication available to us that the descent is a statement about a Frobenius
eigenvalue and not about an asymptotic ratio: the archimedean limit is an accident of the
relative sizes of the characteristic roots, whereas the $p$-adic descent is insensitive to
them. It also means that ``the character of the Ap\'ery limit'' --- which is how we
identified $\chi$ in twelve of the fifteen cases --- cannot be the definition of $\chi$;
for those three it must be read off from the descent itself, or from the modular
parametrisation.

\subsection{Beyond rank two}
\label{ssec:rank3}

The natural next target is the Brown--Zudilin cellular pair for $\zeta(5)$
\cite{BrownZudilin}, where the integral row is a double binomial sum $Q_n$, the second and
third solutions $P_n$ and $\widehat P_n$ satisfy the same order-three recurrence, and the
expected descent has $w=5$. The obstructions are structural and were identified early in
this programme: the local system has rank three, so the descent must be proved for a
three-step filtration at once, with the intermediate weight-three object $\widehat P_n$
appearing in the $P_n$ descent as a cross-term with no $\zeta(3)$ analogue; the summand is
not a perfect square, so the per-carry valuation slack of Remark~\ref{rem:square}
disappears; and the pole to be assembled is of order five rather than three, which roughly
squares the case analysis of Lemma~\ref{lem:V}. The first-row Lucas congruence
$Q_{ap+r}\equiv Q_aQ_r\pmod p$ is proved and formalised (Section~\ref{sec:lean}).

The lesson of Section~\ref{sec:minimal} applies here with force: the right search is not
``which weight-five harmonic monomials span $P_n$'', but ``which weight has an
$n$-difference equal to a single hypergeometric term times a rational function'' --- a
smaller, linear and decidable search, and one insensitive to the large kernel that makes
monomial fitting degenerate. And the top-cell constant, which at weight three is the
classical $5$ of Lemma~\ref{lem:D5}, is computable in advance from the residue at the top
of the support, giving a necessary numerical fingerprint that any candidate weight must
match.

\subsection{Open problems}
\label{ssec:open}

\begin{problem}\label{prob:master}
Prove Conjecture~\ref{conj:master}, in any of its three independent directions: the
multi-digit statement, the sharpening from modulo $p$ to modulo $p^w$, and the removal of
the tameness hypothesis. Even the weight-three case, Conjecture~\ref{conj:master3}, is
open in all three.
\end{problem}

\begin{problem}\label{prob:eps}
Prove Theorem~\ref{thm:brow} by an $\varepsilon$-deformation. Gessel's companion
congruence \cite[Thm.~1.4]{Beukers2211} is structurally the first parameter-derivative of
the first-row congruence; our $b_n$ is the weight-three companion, i.e.\ the third
derivative. Writing $a_n=\ct[\Lambda_\varepsilon^n]$ for a deformed Laurent polynomial,
applying the Dwork--Mellit--Vlasenko ratio congruence \cite{MellitVlasenko} to the
deformed family and taking $\partial_\varepsilon^3|_{\varepsilon=0}$ should land the
weight-three pole and multiply the modulus by $p^3$, since each filtration level costs one
power of $p$ \cite[Prop.~7.7]{BVthree}.
\end{problem}

\begin{problem}\label{prob:lemmaD}
Supply the valuation bound (``a Kummer carry in a wide binomial beats the order-one pole
of a harmonic letter at $\lfloor2k/p\rfloor$'') that would remove the tameness hypothesis
(H4) and put $\alpha$, $\varepsilon$, $s_7$ and $\mathbf E$ inside Theorem~\ref{thm:LB}.
This would also produce the first \emph{proved} instance of the twisted case $e=1$, which
at present has no tame example.
\end{problem}

\begin{problem}\label{prob:seven}
Find harmonic decompositions for the remaining seven families. For $\mathbf C$ this
provably requires letters that its own summand does not supply
(Observation~\ref{obs:negatives}); a constant-term representation \cite{Gorodetsky} is
the natural route. For $\delta$, letters attached to the $3$-section of the summand.
\end{problem}

\begin{problem}\label{prob:zeta2}
Close the $\zeta(2)$-Ap\'ery case of Section~\ref{ssec:obstruction}: either exhibit a
gauge (a shift by an element of the $11$-dimensional kernel of the fitting system) in
which the correction telescopes, or run the argument in a $\Pi\Sigma$ field with nested
sums. Proposition~\ref{prop:obstruction} shows no rational $\times$ harmonic-monomial
ansatz can work.
\end{problem}

\begin{problem}\label{prob:chi-intrinsic}
Give an intrinsic definition of $\chi$ --- one that does not go through the Ap\'ery limit,
and so applies to $\mathbf B$, $\delta$, $\eta$ --- and prove that it is the character of
the quadratic twist of the modular parametrisation attached to the family.
\end{problem}

\section{Machine verification}
\label{sec:lean}

A substantial part of this paper is formalised in Lean~4 and checked by its kernel. This
section reports exactly what is and is not machine-verified. Toolchain:
\texttt{leanprover/lean4:v4.33.0-rc1}, Mathlib revision
\texttt{cd580e54f1a6b46063824e80cec92f64692cbe78}.

\subsection{What is formalised}

\begin{center}\small
\begin{tabular}{@{}ll@{}}
\toprule
result & status\\
\midrule
Theorem~\ref{thm:arow} (first row, Ap\'ery) & formalised, $0$ sorries\\
first row for the Brown--Zudilin double sum $Q_n$ & formalised, $0$ sorries\\
full digit-product form of both first rows & formalised, $0$ sorries\\
Lemma~\ref{lem:K} and Corollary~\ref{cor:K2} & formalised, $0$ sorries\\
Theorem~\ref{thm:LB} (general $\chi$-twisted descent) & formalised, $0$ sorries\\
%% FIXED [MINOR]: the Lean theorem is about $B_{\min}$, not about $b$; Theorem \ref{thm:brow}
%% FIXED is the statement for $b$, and the bridge $B_{\min}=b$ is the unformalised row below.
second row for the minimal form $B_{\min}$ (the $B_{\min}$ analogue of
Theorem~\ref{thm:brow}) & formalised, $0$ sorries\\
second row for Franel & formalised, $0$ sorries\\
two- and three-digit Kummer module (K1--K4) & formalised, $0$ sorries\\
\midrule
$B_{\min}=b$ (Theorem~\ref{thm:minimal}) & \emph{not} formalised (proved on paper)\\
Conjecture~\ref{conj:master} (master form) & not formalised (not proved)\\
\bottomrule
\end{tabular}
\end{center}

\texttt{lake build} is clean; a clean rebuild of all eight modules takes $34$ seconds with
Mathlib prebuilt. A grep for \texttt{sorry}, \texttt{admit} and \texttt{native\_decide}
over the development returns nothing, and \texttt{\#print axioms} on every exported
theorem returns only the three standard Mathlib axioms:
\begin{verbatim}
'ZetaLucas.K_descent'     depends on axioms:
    [propext, Classical.choice, Quot.sound]
'ZetaLucas.mon_descent'   depends on axioms: [ ... same three ... ]
'ZetaLucas.theorem_LB'    depends on axioms: [ ... same three ... ]
'ZetaLucas.bMin_lucas'    depends on axioms: [ ... same three ... ]
'ZetaLucas.bFranel_lucas' depends on axioms: [ ... same three ... ]
'ZetaLucas.apery_lucas'   depends on axioms: [ ... same three ... ]
'ZetaLucas.Q_lucas'       depends on axioms: [ ... same three ... ]
\end{verbatim}
In particular there is no \texttt{sorryAx} and no \texttt{Lean.ofReduceBool}: no
statement rests on a compiled-code evaluation.

\subsection{The statements}

We display the main statements in ASCII-transliterated Lean (Unicode identifiers such as
$\mathbb Q$, $\mathbb N$, $\in$, $\forall$ are written \texttt{Q}, \texttt{N},
\texttt{in}, \texttt{forall}). The predicate \texttt{PCong p x y} is defined as
\texttt{padicNorm p (x - y) < 1}, i.e.\ $x\equiv y\pmod p$ inside $\Zp\subseteq\Q$; it
implies $p$-integrality of both sides as soon as one of them is $p$-integral, so the
statements are not vacuous.

\begin{verbatim}
-- Lemma K
theorem K_descent {chi : N -> Z} (hchi : forall m n, chi (m*n) = chi m * chi n)
    {r : N} (hr : 0 < r) (y : N) :
    PCong p ((p : Q)^r * K chi r y) ((chi p : Q) * K chi r (y / p))

-- Theorem LB (hypotheses H1, H2, H3, H4, H4c, Hw, H5 as in Section 4)
theorem theorem_LB {J : Finset iota} {c : iota -> Q}
    {mon : iota -> List Letter} {S : N -> N -> Z} {w : N} {chip : Z}
    (H1 ...) (H2 ...) (H3 ...) (H4 ...) (H4c ...) (Hw ...) (H5 ...)
    {a r : N} (ha : a < p) (hr : r < p) :
    PCong p ((p : Q)^w * Brow J c mon S (a*p + r))
      ((chip : Q) * Brow J c mon S a * (Arow S r : Q))

-- the second row for the minimal Apery form; bMin n is the sum over
-- k in range (n+1) of (A n k : Q) * (2 * Harm 3 n - Harm 3 k)
theorem bMin_lucas {p : N} [Fact p.Prime] {a r : N}
    (ha : a < p) (hr : r < p) :
    PCong p ((p : Q)^3 * bMin (a*p + r)) (bMin a * (apery r : Q))

-- the second row for Franel
theorem bFranel_lucas {p : N} [Fact p.Prime] (hp2 : p <> 2) {a r : N}
    (ha : a < p) (hr : r < p) :
    PCong p ((p : Q)^2 * bFranel (a*p + r)) (bFranel a * (franel r : Q))
\end{verbatim}

\subsection{Two sharpenings found by formalising}

\begin{enumerate}[leftmargin=1.7em]
\item \texttt{bMin\_lucas} carries \emph{no} hypothesis $p\ge5$ and \emph{no} hypothesis
  $1\le a$. The minimal weight's coefficients are the integers $2,-1$, so the (H4)
  coefficient condition is vacuous and the theorem holds for every prime, including $2$
  and $3$; and $a=0$ is fine, since $B_{\min}(0)=0$ and $p^3B_{\min}(r)$ is $p$-divisible
  for $r<p$. The informal Theorem~\ref{thm:brow} needs $p\ge5$ only because Ap\'ery's
  original weight carries $1/(2m^3)$.
\item \texttt{bFranel\_lucas} needs only $p\ne2$, where the informal analysis had recorded
  $p\ge5$; the coefficient denominators are $4$, and $p=3$ is fine.
\end{enumerate}

\subsection{Design decisions, and why they matter mathematically}

Three changes were made to the informal hypotheses; all three make the theorem easier to
instantiate, not harder.

\begin{enumerate}[leftmargin=1.7em]
\item (H1) is stated as an unconditional congruence
  $S(ap+r,bp+s)\equiv S(a,b)S(r,s)\pmod p$ for \emph{all} $a,b$, together with
  $S(n,k)=0$ for $k>n$; (H2) then disappears. In the carrying regime both sides vanish
  modulo $p$, so the dichotomy is absorbed. The consequence is that \emph{no
  product-region argument occurs anywhere in the formalisation}: the double sum factors
  because a full block $\{0\le k<(a+1)p\}$ splits, and the extension from
  $\{0\le k\le n\}$ to that block is an equality of integers. This retires, for the general
  theorem, the compression that Lemma~\ref{lem:fibre} had to expand by hand.
\item (H3)'s ``surviving set'' is replaced by the side condition
  $S(r,s)\not\equiv0\pmod p$, which is exactly what excludes a borrow in an argument such
  as $n-k$: for Franel it gives $\bin rs\ne0$, hence $s\le r$, hence
  $\lfloor(n-k)/p\rfloor=a-b$. Instances then \emph{derive} the no-borrow condition rather
  than describing $\Sigma_r$.
\item (H5) is stated as ``the product of the letters' characters at $p$ is the same for
  every monomial'' rather than as a count of $\chi$-letters. This is strictly more
  general --- it allows mixed conductors --- and cheaper to check.
\end{enumerate}

The letter system is deliberately concrete: a \texttt{Letter} bundles a completely
multiplicative $\chi:\N\to\Z$, a positive degree, and an argument form
$\N\to\N\to\N$; a monomial is a \texttt{List Letter}, so genuine products such as
$H_kH_{n-k}$ (which Franel needs) are expressible. The sum $\Kchi ry$ is taken over
\texttt{range (y+1)}, so the $m=0$ term is $\chi(0)/0^r=0$ in $\Q$ for $r\ge1$; this makes
the reindexing of the $p$-divisible layer a single bijection $m\mapsto m/p$,
$j\mapsto jp$, and no reindexing of $\texttt{Finset.Icc}$ --- the step that had been
priced as the main schedule risk --- ever arises.

\subsection{Numerical checks inside the build}

Normalisations are pinned in-file, by kernel evaluation where feasible and by
\texttt{\#eval} otherwise:
$\texttt{apery}\mapsto1,5,73,1445,33001,819005$ (A005259, by \texttt{decide});
$Q\mapsto1,21,2989,714549,217515501,76157194521$, confirming the double-sum reading
$Q_1=21$; $\texttt{bMin}\mapsto0,6,\tfrac{351}4,\tfrac{62531}{36},\tfrac{11424695}{288},
\tfrac{35441662103}{36000}$, which is exactly the classical sequence $b_n$;
$\texttt{franel}\mapsto1,2,10,56,346,2252$ (A000172) and
$\texttt{bFranel}\mapsto0,1,4,\tfrac{208}9,\tfrac{1280}9$, matching the Franel recurrence
with $B(0)=0$, $B(1)=1$. An evaluable form of $p$-divisibility sweeps every cell:
both second-row congruences evaluate to \texttt{true} for all $a,r<p$ at $p=5,7$ in the
build, and off-line at $p=11,13$. Sharpness spot-check: at $p=5$, $a=r=1$ the difference
$p^3B_{\min}(ap+r)-B_{\min}(a)\,a_r$ equals $103188530395/32$, of $5$-valuation exactly
$1$ --- the floor is exactly $w$, as the sweeps of Section~\ref{sec:depth} report.

\subsection{Effort, and the economics of the two routes}

The development is $1540$ lines across eight modules
(\texttt{Core}, \texttt{Apery}, \texttt{BrownZudilin}, \texttt{PadicBridge},
\texttt{Letters}, \texttt{TheoremLB}, \texttt{Instances}, \texttt{Kummer}).

The route matters more than the rate. Formalising Theorem~\ref{thm:brow} \emph{via the
proof of Section~\ref{sec:weight3}} --- Lemma~\ref{lem:V}, Corollary~\ref{cor:tailfact},
Lemma~\ref{lem:levelcarry}, the Kummer ledger and the block reindexing of
$\texttt{Finset.Icc}$ --- was costed in advance at $\approx13$ hours $\pm4$, with two
pieces carrying essentially all the schedule risk. Formalising it \emph{via
Theorem~\ref{thm:LB} and the minimal form} cost $\approx2.2$ hours and produced a strictly
stronger result: arbitrary $\chi$, arbitrary weight $w$, all primes, $a=0$ allowed, and a
second sequence for free. The claim made informally --- that the general theorem
``reproves the weight-three theorem in three lines, with no Lemma~V, no tail
factorisation and no Kummer ledger'' --- is thus machine-confirmed, and the cost ratio is
about $6\times$. (Wall-clock figures are reconstructed from the work order rather than
stopwatch readings and should be read as $\pm25\%$; line counts and build-iteration counts
are exact.)

The Kummer module is retained even though the $(\mathrm{LB}_w^\chi)$ route does not use
it: it is precisely the toolkit that the non-tame instances of
Section~\ref{ssec:lemmaD} will need.

\subsection{What is \emph{not} formalised}

\begin{enumerate}[leftmargin=1.7em]
\item \emph{$B_{\min}=b$}. The Lean theorem is about the explicit sum
  $B_{\min}(n)=\sum_kA(n,k)(2\HH3n-\HH3k)$. Its equality with the classical Ap\'ery
  numerator is Theorem~\ref{thm:minimal}, proved on paper here and \emph{not} formalised.
  The gap is therefore a formalisation gap, not a mathematical one, and it is well-scoped:
  Lemma~\ref{lem:conv}'s six conversions, the six certificates of Table~\ref{tab:certs}
  (each an \texttt{Expand}$\to0$ polynomial identity, i.e.\ \texttt{ring} in Lean) and an
  induction on $L$, whose leading coefficient $(n+1)^3$ never vanishes.
\item \emph{$B_{\mathbf A}$ is the second Franel solution}: same status --- it means
  proving that the explicit sum satisfies the Franel recurrence, which is
  Theorem~\ref{thm:franel} here.
\item \emph{The master form}, Conjecture~\ref{conj:master}: not proved, hence not
  formalised. Note that the first-row multi-digit form \emph{is} formalised and is a cheap
  induction; the second row is not, because its ratio form has poles.
\item \emph{A twisted instance.} The $\chi$ machinery of \texttt{theorem\_LB} is proved
  but never exercised at $\chi(p)\ne1$, because the only known twisted sequence
  ($\mathbf E$) is non-tame (Problem~\ref{prob:lemmaD}).
\item \emph{The non-tame instances} and \emph{the seven families with no known
  decomposition}: nothing to formalise until the mathematics exists.
\end{enumerate}

\section{Data availability, and the exact sweeps}
\label{sec:data}

Every computation reported in this paper is exact --- rational arithmetic, or arithmetic
in a finite field with rational reconstruction, or exact polynomial expansion. No floating
point enters any claim, with the single exception of the PSLQ identification of the
Ap\'ery limits (Observation~\ref{obs:limits}), which is a high-precision numerical
identification and is labelled as such.

\subsection{Scripts}

The computations were carried out in the programme's working repository, and the scripts
are archived alongside this paper:

\begin{center}\small
\begin{tabular}{@{}>{\raggedright\arraybackslash}p{5.4cm}>{\raggedright\arraybackslash}p{7.6cm}@{}}
\toprule
location & contents\\
\midrule
\texttt{work/lbw/sporadic.py} & the fifteen sequences, their recurrences and second solutions\\
\texttt{work/lbw/t1b.py} & weights $w$, integrality, Ap\'ery limits (PSLQ)\\
\texttt{work/lbw/t2\_final.py}, \texttt{t2\_diag.py}, \texttt{t2\_twist.py}
   & the master sweeps, the failure diagnosis, the $\chi$-twist\\
\texttt{work/lbw/t3\_fit.py}, \texttt{t3\_fit2.py} & harmonic-decomposition fitting and exact validation\\
\texttt{work/lbw/t4\_proofcheck.py}, \texttt{work/lbw/deep.py} & per-layer checks of the proof of Theorem~\ref{thm:LB}; depth data\\
\texttt{work/minform/step\_a.wl}, \texttt{step\_b.wl}, \texttt{step\_b2.wl}
   & anchors and certificates for Theorem~\ref{thm:minimal}\\
\texttt{work/minform/certs\_stretch.m} & full certificates for Theorems~\ref{thm:franel} and \ref{thm:s10}\\
\texttt{lean/} & the Lean~4 development of Section~\ref{sec:lean} (eight modules)\\
\bottomrule
\end{tabular}
\end{center}

\subsection{The sweeps, with their exact parameters}

\begin{center}\small
\begin{tabular}{@{}>{\raggedright\arraybackslash}p{6.7cm}>{\raggedright\arraybackslash}p{6.3cm}@{}}
\toprule
claim & range, and result\\
\midrule
Theorem~\ref{thm:arow} (first row) & $5\le p\le31$, $n\le320$, single- and multi-digit: $0$ failures\\
Theorem~\ref{thm:brow} (second row, single digit) & $p\in\{5,7,\dots,31\}$, all $(a,r)$: $0$ failures; independently re-implemented for $p\in\{5,7,11,13,17\}$\\
Every lemma and case split of Section~\ref{sec:weight3} & $p\in\{5,7,11,13\}$ (some to $19$), all cells: $0$ failures; $\vp(CC)=3$ branch checked on all $14\,292$ occurrences for $p\le19$; the bound of Lemma~\ref{lem:V} has minimum exactly $1$\\
Conjecture~\ref{conj:master3} (master, $w=3$) & $5\le p\le31$, $n\in[1,320]$: $0$ failures, floor exactly $3$\\
Ratio-form depth on $p$-unit cells & $7\le p\le31$, single-digit: minimum exactly $3$, no cell of depth $\le2$\\
$\vp(b_n)\ge-3\vp(d_n)$ & $5\le p\le31$, $n\le320$: $0$ failures\\
Conjecture~\ref{conj:master} (all fifteen) & every prime $5\le p\le103$, $n\le400$; plus $p=5,7,11$ with $n\le1200$ (five base-$p$ digits): $0$ failures, floor exactly $w$ in every cell\\
Failure of the untwisted law (Obs.~\ref{obs:naive-fail}) & $5\le p\le43$, $n\le300$\\
The repair (Obs.~\ref{obs:repair}) & the seven twisted sequences, $5\le p\le31$ with $\chi(p)=-1$, $n<p^3$ (to $1500$): valuation exactly $w$\\
Ramified prime (Obs.~\ref{obs:ramified}) & $\eta$, $p=5$, $n\le500$: $\min_nv_5(B_n)=0$\\
Small primes (Obs.~\ref{obs:small-primes}) & $p=2,3$, $n\le120$\\
Weights and their sharpness (Obs.~\ref{obs:weights}) & $n\le600$; $\vp(B_n)\ge-w\lfloor\log_pn\rfloor$ over $15\times16\times400$ cells\\
Ap\'ery limits (Obs.~\ref{obs:limits}) & PSLQ at $400$ digits, $33$-constant basis, agreement to $31$--$401$ digits\\
Harmonic decompositions (Table~\ref{tab:decomp}) & fitted mod $2^{61}-1$, rationally reconstructed, validated exactly over $\Q$ on held-out $n=60..72$ ($n=70..82$ for $\mathbf E$)\\
Negative fits (Obs.~\ref{obs:negatives}) & $\mathbf C$: bases of $14$, $44$, $44$ and $152$ elements against $110$ and $220$ exact equations; $\delta$: $65$ elements, $110$ equations; $\mathbf E$: pure bases up to $44$ elements, $8$ arguments\\
Layer-by-layer checks of Theorem~\ref{thm:LB} & $\gamma,\mathbf A,s_{10}$ at $p=5,7,11,13$: $0$ failures; $\mathbf D$ at $p=7,11,13$: $0$ failures, $11$ failing cells at $p=5$; $s_7$: $13$ failing cells at $p=7$, none at $p=5,11,13$\\
Certificates of Table~\ref{tab:certs} & \texttt{Expand}$\to0$; pointwise $n=1..12$, $k=0..n+4$; $B_{\min}-b=0$ exactly for $n\le50$ in two independent implementations\\
Theorems~\ref{thm:franel}, \ref{thm:s10} & certificate identities pointwise $n=1..10$ resp.\ $n=1..9$; $(LB)=0$ numerically for $n\le25$ resp.\ $n\le20$\\
Proposition~\ref{prop:obstruction} & the Gosper linear system has no solution for any $\deg x\le8$\\
\bottomrule
\end{tabular}
\end{center}

\subsection{Provenance}

The results were obtained in a sequence of exact computational campaigns, each recorded in
a dated working memo, and each theorem in this paper was written by transcribing its
statement from the memo that established it and re-checking it against that memo rather
than reconstructing it from memory. Where a memo's own numerical summary disagreed with
the underlying sweep, or where a claim was later superseded, the correction is stated
explicitly in the text: see Section~\ref{ssec:panel} (the ``dip at the reflection centre''
that turned out to be a pole cell), Observation~\ref{obs:false1} and
Observation~\ref{obs:false2} (two false strengthenings), and Section~\ref{ssec:verdict}
(the primes at which the tame proofs actually apply).

\section*{Acknowledgements}

% WORDING FLAGGED FOR RIVER -- please review this sentence before circulation.
This work was carried out in a human--AI collaboration: the conjectures were found by
large exact computational sweeps and the proofs were drafted, refereed and formalised in a
loop between the authors and language-model agents, with every statement re-checked
against its source computation and every proof checked by hand or by machine.


% ---------------------------------------------------------------------------
% BIBLIOGRAPHY
%
% All entries below have been checked against a fetched source (arXiv metadata,
% publisher record, or the text of the cited work itself); see the verification
% memo work/BIBLIO_VERIFY.md for the evidence trail.  No UNVERIFIED- keys remain.
% ---------------------------------------------------------------------------

\begin{thebibliography}{99}

\bibitem{Beukers2211}
F.~Beukers,
\emph{$p$-Linear schemes for sequences modulo $p^r$},
\texttt{arXiv:2211.15240}.

\bibitem{BVone}
F.~Beukers and M.~Vlasenko,
\emph{Dwork crystals I},
\texttt{arXiv:1903.11155}; Int.\ Math.\ Res.\ Not.\ (2021), no.~12, 8807--8844.

\bibitem{BVtwo}
F.~Beukers and M.~Vlasenko,
\emph{Dwork crystals II},
\texttt{arXiv:1907.10390}; Int.\ Math.\ Res.\ Not.\ (2021), 4427--4444.

\bibitem{BVthree}
F.~Beukers and M.~Vlasenko,
\emph{Dwork crystals III: from excellent Frobenius lifts towards
supercongruences},
\texttt{arXiv:2105.14841}; Int.\ Math.\ Res.\ Not.\ (2023).

\bibitem{BVfrob}
F.~Beukers and M.~Vlasenko,
\emph{Frobenius structure and $p$-adic zeta values},
\texttt{arXiv:2302.09603}; Adv.\ Math.\ 480 (2025).

\bibitem{BrownZudilin}
F.~Brown and W.~Zudilin,
\emph{On cellular rational approximations to $\zeta(5)$},
\texttt{arXiv:2210.03391}.

\bibitem{ChamStraub}
M.~Chamberland and A.~Straub,
\emph{Ap\'ery limits: experiments and proofs},
\texttt{arXiv:2011.03400}; Amer.\ Math.\ Monthly 128 (2021), no.~9, 811--824.

\bibitem{Cooper2302}
S.~Cooper,
\emph{Ap\'ery-like sequences defined by four-term recurrence relations},
\texttt{arXiv:2302.00757}.

\bibitem{Delaygue}
\'E.~Delaygue,
\emph{Criterion for the integrality of the Taylor coefficients of mirror maps in several
variables},
\texttt{arXiv:1108.4352}; Adv.\ Math.\ 234 (2013), 414--452.

\bibitem{DRR}
\'E.~Delaygue, T.~Rivoal and J.~Roques,
\emph{On Dwork's $p$-adic formal congruences theorem and hypergeometric mirror maps},
\texttt{arXiv:1309.5902}; Mem.\ Amer.\ Math.\ Soc.\ 246.

\bibitem{Gorodetsky}
O.~Gorodetsky,
\emph{New representations for all sporadic Ap\'ery-like sequences, with applications to
congruences},
\texttt{arXiv:2102.11839}; Exp.\ Math.\ 32 (2023), no.~4, 641--656.

\bibitem{HenningsenStraub}
J.~A.~Henningsen and A.~Straub,
\emph{Generalized Lucas congruences and linear $p$-schemes},
\texttt{arXiv:2111.08641}.

\bibitem{MalikStraub}
A.~Malik and A.~Straub,
\emph{Divisibility properties of sporadic Ap\'ery-like numbers},
\texttt{arXiv:1508.00297}.

\bibitem{MellitVlasenko}
A.~Mellit and M.~Vlasenko,
\emph{Dwork's congruences for the constant terms of powers of a Laurent polynomial},
\texttt{arXiv:1306.5811}.

\bibitem{MirrorNote}
A.~Adolphson and S.~Sperber,
\emph{A note on the integrality of mirror maps},
\texttt{arXiv:2410.04293}.

\bibitem{OsburnSahu}
R.~Osburn and B.~Sahu,
\emph{Supercongruences for Ap\'ery-like numbers},
\texttt{arXiv:0906.3413}; Adv.\ Appl.\ Math.\ 47 (2011), no.~3, 631--638.

\bibitem{Straub1401}
A.~Straub,
\emph{Multivariate Ap\'ery numbers and supercongruences of rational functions},
\texttt{arXiv:1401.0854}; Algebra Number Theory 8 (2014), 1985--2008.

\bibitem{Straub2301}
A.~Straub,
\emph{Gessel--Lucas congruences for sporadic sequences},
\texttt{arXiv:2301.12248}; Monatsh.\ Math.\ (2023).

\bibitem{VMmum}
D.~Vargas-Montoya,
\emph{Maximal unipotent monodromy, congruences \`a la Lucas and algebraic independence},
\texttt{arXiv:2103.15192}.

\bibitem{VMcanon}
D.~Vargas-Montoya,
\emph{$p$-Integrality of canonical coordinates},
\texttt{arXiv:2306.03495}.

\bibitem{Zudilin0202159}
W.~Zudilin,
\emph{An elementary proof of Ap\'ery's theorem},
\texttt{arXiv:math/0202159}.

% ---------------------------------------------------------------------------
% The entries below were flagged as unverified in earlier drafts; each has now
% been checked against a source and the flag removed.
% ---------------------------------------------------------------------------

\bibitem{Gessel1982}
I.~M.~Gessel,
\emph{Some congruences for Ap\'ery numbers},
J.\ Number Theory 14 (1982), no.~3, 362--368.

\bibitem{vdPoorten1979}
A.~van der Poorten,
\emph{A proof that Euler missed\dots\ Ap\'ery's proof of the irrationality of $\zeta(3)$},
Math.\ Intelligencer 1:4 (1978/79), 195--203, \S8.

\bibitem{Nesterenko1996}
Yu.~V.~Nesterenko,
\emph{A few remarks on $\zeta(3)$},
Mat.\ Zametki 59 (1996), no.~6, 865--880;
English transl., Math.\ Notes 59 (1996), 625--636. Lemma~1.

\bibitem{PauleSchneider2003}
P.~Paule and C.~Schneider,
\emph{Computer proofs of a new family of harmonic number identities},
Adv.\ Appl.\ Math.\ 31 (2003), no.~2, 359--378.

\bibitem{Fischler2004}
S.~Fischler,
\emph{Irrationalit\'e de valeurs de z\^eta (d'apr\`es Ap\'ery, Rivoal, \dots)},
S\'em.\ Bourbaki exp.\ 910, Ast\'erisque 294 (2004), 27--62, \S1.2;
\texttt{arXiv:math/0303066}.

\bibitem{Schneider2007}
C.~Schneider,
\emph{Ap\'ery's double sum is plain sailing indeed},
Electron.\ J.\ Combin.\ 14 (2007), \#N5, 3 pp.

\end{thebibliography}


\end{document}
